\part{The marked rank-six construction and its family}
\section{Characteristic-number conventions}\label{pub:V:sec-characteristic-conventions}
\begin{proposition}[The characteristic numbers and middle form]
\hypertarget{tag-034M}{}\label{sec:V-known-inputs}
\StatementRecord{sec:V-known-inputs}{proved}{034M}{}
Let $M$ be a closed connected almost complex manifold of real dimension
eight, with its complex orientation, with $H^2(M;\mathbb Q)=
H^6(M;\mathbb Q)=0$, $b_4(M)=4$, and $\chi(M)=6$.
Then its middle integral pairing modulo torsion is $I_{3,1}$ and
\begin{equation}
 H^4(M;\mathbb Z)/\operatorname{tors}\cong I_{3,1},\qquad
 c_4[M]=6,\qquad c_2^2[M]=2,\qquad
 \sigma(M)=2,\qquad\operatorname{Td}[M]=0.
 \hypertarget{tag-02EQ}{}\label{maps:eq:F.2}
\end{equation}
\end{proposition}
\begin{proof}
The classes $c_1,c_3$ vanish rationally. The index and signature
polynomials of \cite[equations~(1.2)--(1.4), pp.~3--4]{ffEsoleKang}
therefore reduce to
\begin{equation}
 \operatorname{Td}=\frac{3c_2^2-c_4}{720},\qquad
 \chi^1=\frac{3c_2^2-31c_4}{180},\qquad
 \sigma=\frac{3c_2^2+14c_4}{45}.
 \hypertarget{tag-02EF}{}\label{maps:eq:F.1}
\end{equation}
Here the first two numbers are elliptic-symbol indices of the
almost complex tangent bundle; their integrality follows from
\cite[Theorem~24.5.4 and its corollary, p.~184]{HirzebruchTopologicalMethods}.
Substitution of $c_4[M]=6$ gives
$\sigma(M)=16\operatorname{Td}[M]+2$.
The inequality $|\sigma(M)|\le b_4(M)=4$ forces
$\operatorname{Td}[M]=0$ and $\sigma(M)=2$, and then $c_2^2[M]=2$.
Poincar\'e duality makes the middle pairing modulo torsion
unimodular. Its signature excludes an even unimodular form,
whose signature is divisible by eight. The classification of
indefinite odd unimodular forms
\cite[Theorem~2.2.1, p.~61]{ffPetersSterk} gives $I_{3,1}$.
\end{proof}

\begin{remark}[The orientation in the connected-sum criterion]
\hypertarget{tag-PV06}{}\label{pub:V:mandal-convention}
We use the orientation in
\cite[Corollary~1.4, p.~4, and its proof in \S5]{ffMandal}:
$\ell\mathbb{HP}^2\#m(-\mathbb{HP}^2)$ admits an
orientation-compatible almost complex structure exactly for
$\ell=2m+1$ and odd $m$.
The tangent-class conventions are those of
\cite[Lemma~2.1, p.~5, and Theorem~2.4, p.~7]{ffMandal}.
\end{remark}

\section{The monodromy and the marked cusp}\label{pub:V:sec-local}

\begin{definition}[The marked rank-six datum]\hypertarget{tag-PV02}{}\label{pub:V:rank-six-data}\hypertarget{tag-034O}{}\label{sec:fourfold}
\StatementRecord{pub:V:rank-six-data}{proved}{PV02}{}
We use column vectors and the integral cohomology matrices
\[
A=\begin{pmatrix}
1&0&0&1&-1&1\\0&1&-1&-1&1&0\\0&1&0&-1&0&0\\
0&2&-1&-1&0&1\\0&2&0&-1&0&0\\0&1&1&-1&0&0
\end{pmatrix},\quad
B=\begin{pmatrix}
1&0&0&0&0&0\\0&0&-1&0&1&0\\0&0&-1&0&0&1\\
0&0&-1&0&0&0\\0&1&-1&0&0&0\\0&0&-1&1&0&0
\end{pmatrix}.
\]
Put $M_1=A^t$, $M_2=B^t$, $M_0=(M_1M_2)^{-1}$ and
\[
 v_1=(1,3,-2,-1,2,0)^t,\qquad v_2=(-1,0,0,0,0,0)^t.
\]
The affine orders are $(m_1,m_2)=(5,4)$, and $\tau_i(x)=M_ix+v_i/m_i$ on $\mathbb R^6/\mathbb Z^6$.
Positive based meridians satisfy $g_1g_2g_0=1$. The exterior-square basis of $\bigwedge^2\mathbb Z^4$ is $(01,02,03,12,13,23)$.
\end{definition}

\begin{proposition}[the cusp normal form]\hypertarget{tag-02MY}{}\hypertarget{tag-02R7}{}\label{maps:src:F.2.1}
\StatementRecord{maps:src:F.2.1}{proved}{02MY}{the definitions and constructions stated in the proof}


Under the index-one condition of \cref{def:admissible}, the cusp matrix is
conjugate in $\operatorname{GL}_6(\mathbb Z)$ to
\[
 U=\begin{pmatrix}I_3&I_3\\0&I_3\end{pmatrix}.
\]
Its integral centralizer consists exactly of
$\left(\begin{smallmatrix}P&Q\\0&P\end{smallmatrix}\right)$
with $P\in\operatorname{GL}_3(\mathbb Z)$ and
$Q\in\operatorname{Mat}_{3\times3}(\mathbb Z)$.

\end{proposition}


\begin{proof}
\begin{enumerate}[label={$\langle1\rangle$\arabic*.}]
\item Let $N=M_0-I$.  The index-one condition says that the image
is the saturated kernel lattice and that
$\Lambda/\ker N\to\operatorname{im}N$ is an isomorphism.
Choose a basis $a_1,a_2,a_3$ of this image and lifts $b_i$ with
$Nb_i=a_i$.  Saturation shows that the six vectors form an
integral basis.  It gives the displayed matrix.  Commuting a
block matrix with $N$ forces its lower-left block to vanish
and its two diagonal blocks to agree.  Invertibility over
$\mathbb Z$ is then equivalent to unimodularity of that block.

\end{enumerate}
\end{proof}

\begin{theorem}[The rank-six monodromy]\hypertarget{tag-02QI}{}\label{night:P1.1}
\StatementRecord{night:P1.1}{proved}{02QI}{the definitions and constructions stated in the proof}
The finite homology monodromies have orders $5,4$, respectively, and
characteristic polynomials
$(T-1)^2\Phi _5(T)$ and $(T-1)^2(T+1)^2(T^2+1)$.
For $N=M_0-I$ we have $N^2=0$, $\operatorname{rank}N=3$, and
nonzero Smith factors $1,1,1$.  Moreover
\[
 \Lambda^{\langle M_1,M_2\rangle}=0,\qquad
 \Lambda_{\langle M_1,M_2\rangle}=\mathbb Z,
\]
where the quotient homomorphism is $\ell(x)=x_0$.
On the cohomology lattice the corresponding groups are $\mathbb Ze_0$
and $\mathbb Z/2$.  Every nonidentity power of each specified affine
action $\tau_i$ is free.
\end{theorem}
\begin{proof}
\begin{enumerate}[label={$\langle1\rangle$\arabic*.}]
\item Exact multiplication gives
\[
N=\begin{pmatrix}
0&0&0&0&0&0\\0&0&0&1&0&-1\\0&0&0&0&1&1\\
-1&0&0&0&0&0\\1&0&0&0&0&0\\-1&0&0&0&0&0
\end{pmatrix},\qquad
S=\begin{pmatrix}
0&0&0&1&0&0\\0&1&0&0&0&0\\0&0&1&0&0&0\\
-1&0&0&0&1&0\\1&0&0&0&0&1\\-1&0&0&0&0&0
\end{pmatrix}.
\]
The matrix $S$ is unimodular and
$S^{-1}M_0S=\left(\begin{smallmatrix}I_3&I_3\\0&I_3\end{smallmatrix}\right)$.
The columns of $S$ are henceforth denoted
$(a_1,a_2,a_3,b_1,b_2,b_3)$, so that $Nb_j=a_j$.
This is also an integral certificate of the Smith assertion.
The Hermite basis of
$(M_1-I)\Lambda+(M_2-I)\Lambda$ is $(e_1,\ldots,e_5)$;
the stacked matrix of the two fixed-vector equations has rank six.
Transposing the matrices gives a fixed lattice $\mathbb Ze_0$ and
coinvariant Smith factors $(1,1,1,1,1,2)$.  The pairing
$\Lambda^\vee\times\Lambda\longrightarrow\mathbb Z$ identifies the
fixed cohomology vector $e_0$ with the covector $\ell$.

\item We have $M_iv_i=v_i$, $\ell(v_1)=1$, and $\ell(v_2)=-1$.
Thus $\tau_i^k(x)=M_i^kx+kv_i/m_i$, and $\tau_i^{m_i}=1$ on the
torus.  A fixed point of $\tau_i^k$ would give
$k\ell(v_i)/m_i=0$ in $\mathbb R/\mathbb Z$.  For
$0<k<m_i$ this is impossible.  The orders, characteristic polynomials, fixed-space dimensions,
coinvariant Smith factors, primitive cusp image and every affine-power
obstruction are recomputed exactly in
\href{https://github.com/KokunoYumeto/s6-complex-structure-workbench/blob/83d76022bcd442278b760db9ff438afceed6c6ea/pro_returns/S6_SERIES_GRIND_checkpoint_10_20261120T0211Z/code/current_audits/audit_corrections/code/independent_V_connection_spectra_06.py}{\texttt{independent\_V\_connection\_spectra\_06.py}}.
 The displayed matrix $S$ and its
standard-form equality are also verified in
\href{https://github.com/KokunoYumeto/s6-complex-structure-workbench/blob/83d76022bcd442278b760db9ff438afceed6c6ea/pro_returns/S6_SERIES_GRIND_checkpoint_10_20261120T0211Z/code/current_audits/audit_corrections/code/independent_V_markings_06.py}{\texttt{independent\_V\_markings\_06.py}}.

\end{enumerate}
\end{proof}

\begin{theorem}[The Kuhn filling at the cusp]\hypertarget{tag-02R0}{}\label{night:P1.9}
\StatementRecord{night:P1.9}{proved}{02R0}{the definitions and constructions stated in the proof}
The periodic Kuhn fan, in the marking $S$, gives a smooth proper
holomorphic filling of the periods $Z(q)=sI_3+C(q)$ for a sufficiently
small disk.  Its central fibre is reduced with normal crossings.
Modulo the index-one period lattice it has respectively $1,7,12,6$
strata of dimensions $3,2,1,0$, each isomorphic to the corresponding
complex torus $(\mathbb C^*)^d$, and Euler number $6$.
An index-$d$ period sublattice gives $d$ times these stratum counts and
Euler number $6d$.
\end{theorem}
\begin{proof}
The fan consists of the cones over the tetrahedra
\[
 [m,m+e_{\pi_1},m+e_{\pi_1}+e_{\pi_2},m+(1,1,1)],
 \quad m\in\mathbb Z^3,\quad\pi\in S_3.
 \tag{E.4}\hypertarget{tag-029R}{}\label{maps:eq:E.4}
\]
The period-lattice action is
\[
 (x,q)\longmapsto
 \bigl(e^{2\pi iC(q)\lambda}q^\lambda x,q\bigr),
 \qquad\lambda\in\mathbb Z^3,
 \tag{E.5}\hypertarget{tag-029S}{}\label{maps:eq:E.5}
\]

\begin{enumerate}[label={$\langle1\rangle$\arabic*.}]
\item Apply \cite[Theorem~8.7 and Corollary~8.8, arXiv v1, pp.~57--59]{Yang2026}
with $n=3$, $t=q$, $D(q)=C(q)$, $B=I_3$, and the displayed
periodic Kuhn triangulation.  These substitutions give exactly the
action
\[
 (z,q)\longmapsto
 \bigl(q^\lambda\exp(2\pi iC(q)\lambda)z,q\bigr)
\]
and the period order $[I_3\mid sI_3+C(q)]$, where
$q=\exp(2\pi i s)$.  Thus the identity on the dense torus, with
the same integral monomial charts, identifies our quotient with
the cited quotient.  The source theorem supplies its smoothness,
Hausdorffness, properness, and the reduced local equations
$q=x_0\cdots x_k$.  Flatness follows since this nonconstant
function is a nonzero divisor in each smooth chart.
\item For an index-$d$ period sublattice $\Gamma=B\mathbb Z^3$,
the same substitutions give the source's quotient map induced
by $\Gamma\subset\mathbb Z^3$.  It has degree $d$, and
Corollary~8.8 gives $\chi(W)=3!d=6d$ for this marked Kuhn fan.
\end{enumerate}
\end{proof}

\begin{lemma}[Smooth dependence on the bounded cusp correction]
\hypertarget{tag-PV01}{}\label{pub:V:cusp-smooth-dependence}
\StatementRecord{pub:V:cusp-smooth-dependence}{proved}{PV01}{}
Let $C_t(q)$, $0\leq t\leq1$, be a smooth family of holomorphic
$3\times3$ matrix functions on a disc, uniformly bounded on a
smaller closed disc.  The marked Kuhn fillings with periods
$[I_3\mid (\log q/(2\pi i))I_3+C_t(q)]$ are diffeomorphic over
that smaller disc.  The diffeomorphisms preserve the central-fibre
strata and the marked angular and deck lattices.  They can be chosen
on the outer collar to extend a smooth marked trivialization of the
regular torus bundles.
\end{lemma}
\begin{proof}
For the linear deformation $C_t=tC$, \cite[Theorem~8.7, arXiv v1, pp.~57--58]{Yang2026} applies directly to $D(q,s)=sC(q)$ for $|s|<2$.
For a smooth bounded family the same local charts give the following comparison. The toric charts have coordinates $(z_0,z_1,z_2,z_3)$ with
$q=z_0z_1z_2z_3$.  Their changes of coordinates depend smoothly
on $t$ and are toric monomials multiplied by the holomorphic units
$\exp(2\pi iC_t(q)\lambda)$.  They therefore define a smooth
family of manifolds with the same coordinate-stratum charts.
The parameter vector field in each chart lifts $\partial/\partial t$,
fixes $q$, and is tangent to every coordinate stratum.  Differences
of these lifts are vertical fields with the same properties.
A partition of unity gives a global lift with these properties;
on an outer collar take the prescribed lift from the regular marked
torus bundle and use a partition subordinate to that collar and the
interior charts.

For $y=\operatorname{Im}(\log q/(2\pi i))$, uniform boundedness
of the corrections gives, after shrinking the disc,
\[
 \tfrac12\|\lambda\|\leq
 \|(I_3+y^{-1}\operatorname{Im}C_t(q))\lambda\|
 \leq\tfrac32\|\lambda\|,
 \qquad 0\leq t\leq1.
\]
The same fundamental-domain estimate as in the local filling theorem
therefore gives a proper family over the closed disc times $[0,1]$.
The lifted field has a flow for the entire parameter interval.
That flow fixes $q$, preserves each stratum, and has the specified
collar restriction.  Since the parameter interval is contractible
and the lattice marking is constant there, its induced lattice map
is the identity.
\end{proof}

\begin{theorem}[The central cusp fibre and its collapse map]\hypertarget{tag-02QK}{}\label{night:P1.11}\hypertarget{tag-02M6}{}\label{maps:src:E.1.3}
\StatementRecord{night:P1.11}{proved}{02QK}{\SeriesRef{night:P1.9}}
The index-one central cusp fibre $C$ has torsion-free integral homology
with ranks
\begin{equation}
 (b_0,\ldots,b_6)=(1,3,9,9,10,3,1).
 \hypertarget{tag-029T}{}\label{maps:eq:E.6}
\end{equation}
The collapse $\kappa:T^6\to C$ has primitive image in every degree;
it is onto except in degree four, where its cokernel is $\mathbb Z$.
Its image ranks are $(1,3,9,9,9,3,1)$.
In every degree its kernel is the saturation of
$(\bigwedge^kM_0-I)\bigwedge^k\Lambda$.
Dually, the specialization
$\kappa^*:H^k(C;\mathbb Z)\to H^k(T^6;\mathbb Z)$
has image equal to the full monodromy-invariant lattice.
Its kernel is zero except in degree four, where it is $\mathbb Z$.
\end{theorem}
\begin{proof}
\begin{enumerate}[label={$\langle1\rangle$\arabic*.}]
\item Apply \cref{pub:V:cusp-smooth-dependence} to $C_t(q)=tC(q)$.
The comparison preserves the central strata, the integral marking,
and the real oriented logarithmic boundary; it therefore intertwines
the collapse maps. We calculate with $C(q)=0$.

\item Here is a concrete topological model.  Put $G=4I_3-\mathbf1\mathbf1^t$.
Its positive eigenvalues are $1,4,4$.  The Voronoi polytope in dual
coordinates is cut out by
\[
|v^ty|\leq \tfrac12v^tGv
 \quad (0\ne v\in\{0,1\}^3).
\]
The value on a binary vector with $k$ nonzero entries is $k(4-k)$.
These fourteen primitive facet normals give the star of a vertex in
the Kuhn triangulation.  One checks its Delaunay tetrahedra by the
successive chains
$0,e_i,e_i+e_j,e_1+e_2+e_3$ for the six permutations $(i,j,k)$;
solving their three perpendicular-bisector equations and the fourteen
inequalities verifies precisely their common faces.  The translates
of this Voronoi polytope by $G\mathbb Z^3$ tile $\mathbb R^3$.

\item For a simplex $\sigma$ with vertices $v_0,\ldots,v_r$, put
$L_\sigma=\sum_i\mathbb Z(v_i-v_0)$ and
$D_\sigma=\bigcap_{v\in\sigma}P_v$, where $P_v=P+Gv$.
The component torus has lattice $A=\mathbb Z^3$; the common-stratum
marking is the explicit map
\[
 (A\oplus\mathbb Z)/\sum_i\mathbb Z(v_i,1)
 \longrightarrow A/L_\sigma,\qquad
 [(a,k)]\longmapsto a-kv_0.
\]
The Kuhn unimodularity makes $L_\sigma$ primitive.

We choose compatible moment maps as follows.  The polytope $2P_v$
is integral, since its vertex facet normals are unimodular and its
facet constants are integral.  Embed the component $X_v$ using the
characters $\chi^{-m}$ for
$m\in\mathcal A_v=2P_v\cap A^\vee$, and define
\[
 \mu_v([Z_m])=
 \frac{\sum_{m\in\mathcal A_v}m|Z_m|^2}
 {2\sum_{m\in\mathcal A_v}|Z_m|^2}.
\]
The sign is specified: the ray $w-v$ tends to the face on which
$\langle w-v,-m\rangle$ is minimal, hence to the upper common
Voronoi face of $P_v$ and $P_w$.  Thus the character polytope is
$-2P_v$, and $\mu_v$ is minus one-half its usual moment map.
The monomial moment theorem
\cite[equations (12.2.3)--(12.2.4) and Theorem 12.2.5, p.~571]{CoxLittleSchenck2011}
identifies the nonnegative toric part with $P_v$.
For the squared weights here, the intervening map on that part is
$[a_m]\mapsto[a_m^2]$, with inverse $[b_m]\mapsto[\sqrt{b_m}]$:
both preserve the homogeneous binomial relations and positive common
projective rescaling.  Write $s_v:P_v\to X_v$ for the resulting
continuous positive section.

On a shared face put
$\mathcal A_\sigma=2D_\sigma\cap A^\vee$ and choose
$m_0\in\mathcal A_\sigma$.
For $m\in\mathcal A_\sigma$, the difference $m-m_0$ annihilates
$L_\sigma$.  The ratio $Z_m/Z_{m_0}$ is consequently the restriction
of the same character
\[
 \bigl(-(m-m_0),\langle m-m_0,v\rangle\bigr)
       \in(A\oplus\mathbb Z)^\vee
\]
for every vertex $v$ of $\sigma$.  At $C=0$ all deck unit multipliers are one, and the toric gluing
identifies these character functions exactly.  On either component
the restricted moment map is therefore the identical expression
\[
 \frac{\sum_{m\in\mathcal A_\sigma}
       m|\chi^{-(m-m_0)}|^2}
 {2\sum_{m\in\mathcal A_\sigma}|\chi^{-(m-m_0)}|^2}.
\]
The equality persists on every smaller face by restricting the index
set.  Positive character values also agree, so the sections $s_v$
glue on the entire shared orbit closure.  The angular maps are the
displayed quotient homomorphisms $A\to A/L_\sigma$.

The period deck map is $(a,k)\mapsto(a+k\nu,k)$, sending $v$ to
$v+\nu$.  Its component marking is unchanged because
$a+k\nu-k(v+\nu)=a-kv$.
Moreover $\mathcal A_{v+\nu}=\mathcal A_v+2G\nu$; the common
character factor cancels in projective ratios.  Hence
\[
 \mu_{v+\nu}(T_\nu p)=\mu_v(p)+G\nu,
 \qquad s_{v+\nu}(y+G\nu)=T_\nu s_v(y).
\]
These give an explicit well-defined map
\[
 \begin{gathered}
 \mathcal Q=
 \bigl((\mathbb R^3/G\mathbb Z^3)\times
       (\mathbb R^3/\mathbb Z^3)\bigr)/\!\sim
 \longrightarrow C,\\
 [y,\theta]\longmapsto e^{2\pi i\theta}\,s_v(y)
 \quad(y\in P_v).
 \end{gathered}
\]
where over $\operatorname{relint}D_\sigma$ the relation collapses
the subtorus $L_\sigma\otimes(\mathbb R/\mathbb Z)$.
On each stratum the inverse takes its moment value and its angular
character class in $T_A/T_{L_\sigma}$; the latter is obtained by
$z\mapsto z/|z|$ on its nonzero characters.
The map is a continuous bijection, and its source is compact and
$C$ Hausdorff, so it is a homeomorphism.
The angular identity $[y,0]$ gives the continuous section, with
the same face and deck compatibilities just proved.

\item The nearby-fibre collapse is the quotient map in this model.
Give the semistable family its divisorial logarithmic structure.
The local equations $q=x_0x_1x_2x_3$ and their face restrictions
make it fine, saturated, log smooth and vertical over the standard
logarithmic disc.  The family is proper by \cref{night:P1.9}.
In its logarithmic central fibre, retain all arguments of the $x_j$
with their sum fixed by $\arg q$.  The integral angular lattice over
a face with vertices $v_0,\ldots,v_r$ is the kernel of the height
map on $\mathbb Z^3\oplus\mathbb Z$.  Forgetting the arguments of
the zero coordinates induces exactly
\[
 \mathbb Z^3\longrightarrow
 (\mathbb Z^3\oplus\mathbb Z)/
       \sum_i\mathbb Z(v_i,1)
 \ \cong\ \mathbb Z^3/L_\sigma.
\]
The last isomorphism sends the class of $(a,k)$ to
$a-kv_0\pmod {L_\sigma}$; its inverse sends $a$ to $(a,0)$.

The polar maps agree under the integral monomial overlaps.
At $C(q)=0$ the deck multipliers have positive real values, so the
logarithmic fibre at angle one is the product of the base and angular
three-tori in item~3, with their given markings.

The mapping-cylinder comparison for a proper vertical log-smooth
family identifies this logarithmic forgetful map with the nearby-fibre collapse, compatibly with the central strata
\cite[formula (3-3-1), p.~478; Theorem 4.1.1 and Lemmas
4.1.2--4.1.3, pp.~480--482]{ffAchingerOgus20}.
Its angular cycles are $a_j$ and its base cycles are $b_j$.
In these column-basis conventions monodromy fixes each $a_j$
and sends $b_j$ to $b_j+a_j$, as specified by $Nb_j=a_j$.
The stratified deformation in item~1 transports this map
back to the original bounded period matrix.

\item Filter $\mathcal Q$ by the dimension $d$ of its dual base cells.
Over the interior of the dual cell of a $p$-simplex, $d=3-p$ and the
angular quotient has lattice
$Q_\sigma=\mathbb Z^3/L_\sigma$ of rank $d$.
Excision for the cell and its boundary gives the relative chain
groups of the resulting homology spectral sequence:
\[
E^1_{d,q}=\bigoplus_{\dim\sigma=3-d}\bigwedge^qQ_\sigma.
\]
Orient dual cells by the orientation of the base and the ordered
simplex.  The differential is the signed quotient map for each
codimension-one coface.  More explicitly, with integral projection
matrices $Q_\sigma$ and chosen right inverses $S_\sigma$,
\[
D_{p,q}[\tau,\sigma]=
 \sum_{\sigma=\partial_j\tau\ {\rm mod}\ \mathbb Z^3}
 (-1)^j\bigwedge^q(Q_\tau S_\sigma).
\tag{1.11.1}
\]
All $L_\sigma$ are primitive because the tetrahedra are unimodular.
Thus these projections and right inverses are integral.  Formula
(1.11.1) specifies every integer entry, including multiple incidences
under translation identifications.  The program \href{https://github.com/KokunoYumeto/s6-complex-structure-workbench/blob/83d76022bcd442278b760db9ff438afceed6c6ea/pro_returns/S6_SERIES_GRIND_checkpoint_10_20261120T0211Z/code/current_audits/audit_corrections/code/reconstruct_V_cusp_06.py}{\texttt{reconstruct\_V\_cusp\_06.py}} enumerates all $26$ simplex classes and $74$ codimension-one incidences and computes these matrices exactly.

\item The nonzero $E^2_{d,q}$ ranks are
\[
\begin{array}{c|rrrr}
q\backslash d&0&1&2&3\\\hline
0&1&3&3&1\\
1&0&6&8&3\\
2&0&0&7&3\\
3&0&0&0&1
\end{array}
\]
and every group is free. The reconstructed incoming Smith matrices
have only unit nonzero factors and reproduce every entry of this table.
Support leaves only the possible higher differential
$d_2:E^2_{3,0}\to E^2_{1,1}$.
The angular identity section splits the base projection as a map of
filtered spaces, so the base fundamental class survives.  This kills
that differential.  All subsequent differentials have zero source or
target.  The filtration extensions are extensions of free abelian
groups and therefore split, proving the asserted integral homology.

\item The product-to-quotient map induces on $E^1$ the exterior powers of
the quotient projections.  The program \href{https://github.com/KokunoYumeto/s6-complex-structure-workbench/blob/83d76022bcd442278b760db9ff438afceed6c6ea/pro_returns/S6_SERIES_GRIND_checkpoint_10_20261120T0211Z/code/current_audits/audit_corrections/code/reconstruct_V_cusp_06.py}{\texttt{reconstruct\_V\_cusp\_06.py}} constructs the induced maps
on integral cycle and quotient bases from the face projections and the
Poincar\'e-dual Alexander--Whitney formula in step 9. Their images are primitive, and the only nonzero graded cokernel
is $\mathbb Z$ at $(d,q)=(2,2)$.  The summed image ranks are
\[
             r_k=(1,3,9,9,9,3,1).
\]
The same exact exterior-power calculation gives these ranks for
$\ker(\bigwedge^kM_0-I)$.
The collapse satisfies
$\kappa_*(\bigwedge^kM_0-I)=0$, because a monodromy circuit
has the same retraction to the central fibre in the mapping-cylinder
comparison.  Therefore its rank is at most $r_k$.
A filtered map has rank at least the sum of the ranks of its
graded maps; hence the rank is exactly $r_k$.

Choose integral bases compatible with the source and target
filtrations; their graded groups are free.  In each graded block
choose the unit Smith pivots of the computed map.  The corresponding
$r_k$ rows and columns of the filtered map form a block
triangular submatrix with unimodular diagonal pivot blocks.
Its determinant is $\pm1$.  Since the full map has rank $r_k$,
these unit pivots prove that its image is primitive.
Thus the map is onto in every degree except four, where its
cokernel is the free group of rank $10-9=1$.
This rank argument also rules out additional lower-filtration images
of source classes in the kernels of the graded maps.

\item We represent a base cycle $b_J$ by
Poincar\'e duality from the Alexander--Whitney cup cocycle on the
complementary coordinates.  On an ordered $p$-simplex its coefficient
is
\[
\operatorname{sgn}(J,J^c)
 \prod_{r=1}^p(v_r-v_{r-1})_{(J^c)_r}.
\]
Tensoring with an angular exterior vector and applying (1.11.1)
gives the marked matrices. Their source basis is the ordered exterior
basis $(a_1,a_2,a_3,b_1,b_2,b_3)$; the output gives the integral
target homology projections and all Smith basis maps.
They annihilate $\bigwedge^kM_0-I$ exactly.  The collapse also
annihilates that operator, since turning the smoothing parameter
does not change its image in the central fibre.  The computed ranks
equal the ranks of the monodromy coinvariants; primitivity therefore
identifies both kernels with the saturation of its image.
Consequently these are the matrices of the collapse in suitable integral target bases.

\item Every homology group just computed is free, so the universal
coefficient maps identify cohomology with the integral duals of homology.
The transpose of the collapse has image
\[
 \operatorname{im}\kappa^*
 =\operatorname{Ann}(\ker\kappa_*)
 =\operatorname{Ann}\bigl((\bigwedge^kM_0-I)\bigwedge^k\Lambda\bigr)
 =\ker\bigl((\bigwedge^kM_0)^t-I\bigr).
\]
The first equality uses the primitive image of $\kappa_*$: an integral
functional on that image extends to the free central-fibre group.
Passing to saturation preserves the annihilator. The last group is
the full integral cohomology invariant lattice; inverse transpose
has the same fixed subgroup. The kernel of $\kappa^*$ is the
integral dual of $\operatorname{coker}\kappa_*$, giving precisely the
stated degree-four kernel.
\end{enumerate}
\end{proof}

\begin{remark}[The branch resolution of the Kuhn fibre]
\hypertarget{tag-PV08}{}\label{pub:V:cusp-branch-resolution}
The branch resolution and the moment-cell complex have the following
augmentation. Over a dual cell $D_\eta$ of the Kuhn decomposition,
the angular lattice is $Q_\eta=\mathbb Z^3/L_\eta$ and the branch
indices form the simplex on the vertices of $\eta$. Tensor its
augmented integral simplex complex with $\bigwedge^qQ_\eta$ and
sum over dual cells. The degree-zero augmentation adds branch
coefficients; the higher branch terms map to zero. The incidence
maps $Q_\eta\to Q_\tau$ for cofaces $\tau\supset\eta$ commute
with this augmentation. Inserting the least vertex contracts each
augmented simplex complex. Thus the augmentation is a
quasi-isomorphism to the moment-cell complex of
\cref{night:P1.11}. Taking cohomology of the compact toric branch
intersections first gives the normalization table below.

The normalization is the three-dimensional permutohedral toric variety.
Its paired boundary surfaces consist of four copies of $\mathrm{dP}_6$
and three copies of $\mathbb P^1\times\mathbb P^1$.  There are twelve
normalized triple curves and six quadruple points.  The normalization
complex for the constant sheaf is the alternating branch complex.
Local exactness follows from the augmented simplex on the branches
of $z_0z_1z_2z_3=0$.  Its nonzero rows after taking cohomology are
\[
\begin{array}{c|rrrr|rr}
q& p=0&p=1&p=2&p=3&\operatorname{rank}d^0&
 \operatorname{rank}d^1\ \text{or }d^2\\ \hline
0&1&7&12&6&0&4,5\\
2&11&22&12&0&5&9\\
4&11&7&0&0&4&0\\
6&1&0&0&0&0&0
\end{array}
\tag{E.7}\hypertarget{tag-029U}{}\label{maps:eq:E.7}
\]
Every nonzero Smith invariant in these matrices is one.


The integral matrices are as follows.
The fourteen rays are $\pm v$, with
$v\in\{0,1\}^3\setminus\{0\}$.  The twenty-four maximal cones
are the translated Kuhn tetrahedra containing zero.  This enumeration
requires exactly the eight cubes with lower vertex in $\{-1,0\}^3$;
no omitted cube can contain zero.  The Picard group of the normalization
is the group of ray coefficients modulo the three character relations.
We set the coefficients at $e_1,e_2,e_3$ to zero, giving an integral
basis of rank eleven.  On a boundary surface $D_v$, choose adjacent
rays $u_1,u_2$ for which $(v,u_1,u_2)$ is unimodular and set the two
remaining character coefficients to zero.  Restriction is evaluated
by subtracting the unique integral linear function with the prescribed
values at these three rays.  The paired restriction is
$\operatorname{res}_{D_v}-\operatorname{res}_{D_{-v}}$, with the
second adjacent ray corresponding to $u-v$.


For a triangle $(0,v,u)$ with $0<v<u$ coordinatewise, the next map
is the alternating sum of the degrees on the three boundary curves
of $D_{u-v},D_u,D_v$.  In a surface fan, if the neighbours of a ray
$w$ are $h_1,h_2$ and $h_1+h_2=bw$ modulo $v$, that degree is
$a_{h_1}+a_{h_2}-ba_w$.  In degree four, the restriction map is dual
to the seven divisor differences $[D_v]-[D_{-v}]$.  In degree zero
the matrices are the cellular incidence matrices of the one-vertex
Kuhn triangulation of $T^3$.  These rules define \eqref{maps:eq:E.7} over
$\mathbb Z$, and both compositions are zero.
Its exact Smith reductions give the ranks and primitivity just stated.


\end{remark}

\begin{lemma}[The marked smooth gluing]
\hypertarget{tag-PV11}{}\label{pub:V:smooth-gluing}
\StatementRecord{pub:V:smooth-gluing}{proved}{PV11}{}
The datum of \cref{pub:V:rank-six-data}, with its two affine caps,
the primitive Kuhn cusp in the marking $S$, and zero additional
translation clutch, gives a closed connected oriented smooth
eight-manifold $M$. It is simply connected, has Euler number six,
and admits an almost complex structure with this orientation.
\end{lemma}
\begin{proof}
\begin{enumerate}[label=\arabic*.,leftmargin=*]
\item Use the matrices and affine vectors of \cref{pub:V:rank-six-data}.
Theorem~\ref{night:P1.1} gives their orders, the primitive cusp
marking, the coinvariant homomorphism $\ell(x)=x_0$, and freeness of
every nonidentity power of the affine actions. Consequently
\[
 (T^6\times\Delta_u)/
 \langle(x,u)\mapsto(M_ix+v_i/m_i,\zeta_{m_i}u)\rangle,
 \qquad (m_1,m_2)=(5,4),
\]
are smooth oriented finite caps. On the universal angular cover
$u=r e^{2\pi ia}$ their marked collar maps are
\[
 \alpha_i(x,u)=x-av_i,
 \qquad
 \alpha_i(M_ix+v_i/m_i,\zeta_{m_i}u)=M_i\alpha_i(x,u).
\]
At the cusp we take the index-one Kuhn cap of
Theorem~\ref{night:P1.9}, in the marking of
Proposition~\ref{maps:src:F.2.1}. The three displayed boundary bundles
therefore glue to the marked torus bundle over the intervening pair
of pants, with zero additional clutch translation. Every boundary
is paired. The result is a closed connected smooth eight-manifold,
oriented by the cusp complex structure and the continued fibre/base
orientation.

\item The small cusp cap retracts onto its central fibre $C$ through
the mapping-cylinder homeomorphism
\cite[Theorem 4.1.1 and Lemmas 4.1.2--4.1.3, pp.~480--482]{ffAchingerOgus20}.
Indeed its divisorial log chart is
$\mathbb N\to\mathbb N^r$, $1\mapsto(1,\ldots,1)$; it is log smooth,
saturated and vertical, and the cap is proper by
Theorem~\ref{night:P1.9}. In the resulting cylinder the homotopy
$[x,r]\mapsto[x,(1-t)r]$ fixes $C$ and retracts onto it.

Before the period quotient, the components of $C$ and their
nonempty intersections are smooth complete toric strata. Each is
simply connected: a loop moves into its dense torus by general
position, and the one-parameter loops along the rays of a smooth
maximal cone bound the coordinate discs; those rays are an integral
basis. Van Kampen sends component-crossing paths to edges and triple
intersections to triangles of the Kuhn nerve in $\mathbb R^3$.
Thus the covering core has trivial fundamental group. Its period
quotient has deck group $\mathbb Z^3$, so $\pi_1(C)=\mathbb Z^3$.
The nearby-torus map is the quotient of $\Lambda$ by the saturated
image of $M_0-I$, by the $H_1$ collapse in
Theorem~\ref{night:P1.11}. It is onto, and the cusp meridian dies
because its zero section extends.

We also record the invariant normal closure needed below. Let
$K$ be the smallest $\langle M_1,M_2\rangle$-invariant subgroup
containing $(M_0-I)\Lambda$, and let
$J=(M_1-I)\Lambda+(M_2-I)\Lambda=\ker\ell$.
The equality $\ell M_i=\ell$ gives $K\subseteq\ker\ell$.
On the abelian quotient $\Lambda/K$, the induced maps $a,b$
satisfy $ab=1$, $a^5=1$ and $b^4=1$. Thus $a^4=a^5=1$, so
$a=b=1$ and $J\subseteq K$. Consequently $K=J=\ker\ell$ as
integral subgroups.


Let $a_i$ be the two finite peripheral lifts. The universal covering
deck transformation of the $i$th cap has $m_i$th power equal to
translation by $v_i$, since $M_iv_i=v_i$. The collar map in
Step~1 sends $a_i$ to that transformation. Van Kampen therefore gives
\[
 a_i\lambda a_i^{-1}=M_i\lambda,
 \qquad a_1^5=v_1,\quad a_2^4=v_2,\quad a_1a_2=1.
\]
Write $a_1=t$ and $a_2=t^{-1}$. Then
$t=t^5t^{-4}=v_1+v_2$ lies in the abelian fibre subgroup. Conjugation
is consequently trivial, and that subgroup is a quotient of
$\Lambda_{\langle M_1,M_2\rangle}=\mathbb Z$, with coordinate $\ell$.
Now $\ell(v_1+v_2)=0$, so $t=0$, while $t^5=v_1$ kills its generator
$\ell(v_1)=1$. Every generator vanishes, proving $\pi_1=1$.
Equivalently, after this abelian reduction, the relations on
$(\ell,a_1,a_2)$ have matrix
\[
 \begin{pmatrix}-1&5&0\\1&0&4\\0&1&1\end{pmatrix},
 \qquad\det=-1.
\]

\item Euler additivity over the base strata gives $\chi(M)=6$:
the regular tori and finite free torus quotients have Euler number
zero, and \cref{night:P1.9} gives six for the cusp.
\item Each finite monodromy admits an invariant real linear complex
structure: its $1$ and $-1$ real eigenspaces have even dimension, and
each remaining conjugate eigenpair gives a real two-plane. Choose
a complex structure on each summand. Reversing it on one fixed
two-plane selects the required orientation. These structures commute
with $M_i$, so their direct sums with the root-disc structure descend
to the affine caps. The cusp cap already has its complex
structure. Their vertical complex structures on the boundary collars
are translation-invariant in the marked torus fibres.

The space $\mathcal J=\operatorname{GL}_6^+(\mathbb R)/
\operatorname{GL}_3(\mathbb C)$ of allowable vertical structures
retracts to $\operatorname{SO}(6)/\operatorname{U}(3)$ by the explicit
conjugation
\[
 J\longmapsto P_tJP_t^{-1},
 \qquad P_t=\bigl((1-t)I+t(I+J^tJ)\bigr)^{1/2}.
\]
At $t=1$ the structure is orthogonal, since $J$ preserves the metric
$I+J^tJ$; an orthogonal $J$ is fixed throughout this homotopy.
The inclusion $\operatorname{U}(3)\to\operatorname{SO}(6)$ sends
$\operatorname{diag}(e^{i\theta},1,1)$ to one-plane rotation and
induces the surjection $\mathbb Z\to\mathbb Z/2$ on fundamental
groups. Indeed the determinant fibration and
$\operatorname{SU}(2)=S^3\to\operatorname{SU}(3)\to S^5$ give the
first group; the quaternion cover of $\operatorname{SO}(3)$ and
the sphere fibrations for $\operatorname{SO}(n)$ give the second,
with a full rotation lifting from $1$ to $-1$. The homogeneous-space
exact sequence therefore gives $\pi_0\mathcal J=\pi_1\mathcal J=0$.

Over the two-dimensional pair of pants the prescribed boundary
sections of the associated $\mathcal J$-bundle extend over relative
one-cells by connectedness and over relative two-cells by simple
connectivity. Smooth approximation relative to the collars gives
a smooth vertical structure. Horizontal splittings form an affine
space, so a partition of unity extends the prescribed complex
splittings from those collars. Combining the vertical structure with
the pulled-back base structure and gluing the caps gives the required
almost complex structure on $M$.

\end{enumerate}
\end{proof}

\section{The genus-two pencil and the lifted period planes}



\begin{theorem}[The genus-two pencil and its Hodge bundle]\hypertarget{tag-02QV}{}\label{night:P1.4}
\StatementRecord{night:P1.4}{proved}{02QV}{the definitions and constructions stated in the proof}
The smooth projective models of the pencil
\[
 C_\lambda:\quad y^2=x^5+\lambda^3x+\lambda^4
 \tag{F.7}\hypertarget{tag-02FB}{}\label{maps:eq:F.7}
\]
have one nodal fibre away from $0,\infty$, at $\lambda_*=-256/3125$.
The minimal good root orders at $0,\infty$ are $8,5$.
On the root curve with reduced root divisors $D_0=[0]/8$ and $D_\infty=[\infty]/5$, the Hodge bundle is
\[
 \mathcal E=\mathcal T^3\oplus\mathcal T,\qquad
 \mathcal T=\mathcal O_{\mathcal B'}(2D_\infty-3D_0),\qquad
 \deg\mathcal T=\frac1{40}.
 \tag{F.8}\hypertarget{tag-02FC}{}\label{maps:eq:F.8}
\]
Thus $T=\mathcal T=\mathcal O((2/5)[\infty]-(3/8)[0])$, and $\ell_0=T^3$ is generated meromorphically by $dx/y$.
The extension at $\lambda_*$ is the dualizing Hodge extension.
\end{theorem}
\begin{proof}
\begin{enumerate}[label={$\langle1\rangle$\arabic*.}]
\item The discriminant of $x^5+\lambda^3x+\lambda^4$ is
$\lambda^{15}(3125\lambda+256)$.
At $\lambda_*$ the double root is $64/625$;
the second $x$-derivative and first $\lambda$-derivative there are
nonzero.  The remaining roots are simple.  The odd-degree
hyperelliptic compactification is smooth at its branch point at
infinity.  This gives one node, a smooth total surface at that node,
and genus-one normalization.

\item At zero set $\lambda=u^8$, $x=u^6X$, $y=u^{15}Y$; at infinity set
$\lambda=u^{-5}$, $x=u^{-4}X$, $y=u^{-10}Y$.
The resulting smooth equations and differentials are
\[
\begin{array}{c|c|c|c}
 &\text{equation}&dx/y&x\,dx/y\\\hline
0&Y^2=X^5+X+u^2&u^{-9}dX/Y&u^{-3}X\,dX/Y\\
\infty&Y^2=X^5+uX+1&u^6dX/Y&u^2X\,dX/Y.
\end{array}
\]
Their central polynomials are square-free.  The deck maps at zero and
infinity are $(X,Y)\mapsto(iX,\zeta_8Y)$ and
$(X,Y)\mapsto(\zeta_5^{-1}X,Y)$, respectively.
The eigenvalues on $H^{1,0}$ are $(\zeta_8,\zeta_8^3)$ and
$(\zeta_5^{-1},\zeta_5^{-2})$.  Including conjugates gives characteristic
polynomials $T^4+1$ and $\Phi_5$, proving minimality of the root orders.
The fourth power at zero is the hyperelliptic involution, acting as
$-I$ on $H^1$ and as $I$ on its exterior square.

\item The two meromorphic differential generators have orbifold divisors
$(6/5)[\infty]-(9/8)[0]$ and
$(2/5)[\infty]-(3/8)[0]$.
At the node they are independent sections of the dualizing Hodge
extension: one has opposite nonzero residues at the two points of the
normalization, and the linear combination vanishing at the node is a
regular nonzero elliptic differential.  Thus there is no additional
divisor at $\lambda_*$.  These computations prove the asserted direct
sum and its degrees.  The discriminant, the nodal root and its transverse derivatives,
both exact root substitutions, square-freeness of the central polynomials,
and the line degree are recomputed in
\href{https://github.com/KokunoYumeto/s6-complex-structure-workbench/blob/83d76022bcd442278b760db9ff438afceed6c6ea/pro_returns/S6_SERIES_GRIND_checkpoint_10_20261120T0211Z/code/current_audits/audit_corrections/code/independent_V_markings_06.py}{\texttt{independent\_V\_markings\_06.py}}, with its exact JSON output.

\end{enumerate}
\end{proof}









\begin{theorem}[The rational rank-five monodromy]\hypertarget{tag-02O7}{}\label{maps:src:F.3.5}
\StatementRecord{maps:src:F.3.5}{proved}{02O7}{\SeriesRef{maps:src:F.2.5}}


For $U=\mathbb P^1\setminus\{0,\infty,\lambda_*\}$, let
$\mathcal H_{\rm dR}$ be the polarized rational cohomology system
of the pencil in \cref{night:P1.4}. Its primitive exterior square
is rationally isomorphic to the rank-five quotient of the system
with matrices $A,B$ in \cref{pub:V:rank-six-data} by its invariant
first-coordinate line. The identification uses the explicit
companion and quotient maps below and the stated based meridians.
\end{theorem}


\begin{proof}
\begin{enumerate}[label={$\langle1\rangle$\arabic*.}]
\item By \cref{night:P1.4}, the cusp is an ordinary node with
genus-one normalization. Cutting and regluing an annular neighborhood
gives a Dehn twist in the primitive nonseparating vanishing circle.
Its logarithm $N_H$ has square zero and rank one.
The finite characteristic polynomials, from the same theorem, are
$f(T)=T^4+T^3+T^2+T+1$ at infinity and $g(T)=T^4+1$ at zero.

\item We identify the representation.
Choose positive peripheral generators $T_1,T_2$ at infinity and
zero, in that order, with $T_1T_2$ inverse to nodal monodromy.
Set $R=T_1^{-1}$ and $S=T_2$.  Then $R-S$ has rank one and their
characteristic polynomials are $f,g$.  The companion presentation
is Levelt's theorem in \cite[Theorem 3.5, pp.~331--332;
Corollary 3.6, p.~332]{BeukersHeckman1989}.  Its generators are
related to the present pair by
\[
 h_\infty=R=T_1^{-1},\qquad h_0=S^{-1}=T_2^{-1},\qquad
 h_1=R^{-1}S=T_1T_2,\qquad h_\infty h_1h_0=I.
\]
The following calculation supplies the rational change
of basis for this presentation.
If $u$ spans $\operatorname{im}(R-S)$, it is cyclic for $R$:
the span of its iterates is stable under both $R$ and $S$, and
the irreducibility of $f$ makes that span all of $\mathbb Q^4$.
In the rational basis $u,Ru,R^2u,R^3u$, the matrix of $R$ is
the companion matrix of $f$, and $S=R-u\varphi$ for a covector
$\varphi$.  The determinant identity
\[
 g(T)=f(T)+\varphi\operatorname{adj}(TI-R)u
 \tag{F.9}\hypertarget{tag-02FD}{}\label{maps:eq:F.9}
\]
determines $\varphi$ uniquely: the four coordinates of
$\operatorname{adj}(TI-R)u$ are a basis of the polynomials of
degree at most three, as is seen from their successively distinct
leading degrees in the cyclic basis.  It follows that any two
such pairs are simultaneously rationally conjugate.  Apply the
lemma to the explicit companion pair
\[
 C_f=\begin{pmatrix}
 0&0&0&-1\\1&0&0&-1\\0&1&0&-1\\0&0&1&-1
 \end{pmatrix},\qquad
 C_g=\begin{pmatrix}
 0&0&0&-1\\1&0&0&0\\0&1&0&0\\0&0&1&0
 \end{pmatrix}.
 \tag{F.10}\hypertarget{tag-02EG}{}\label{maps:eq:F.10}
\]
Their difference has rank one and their characteristic polynomials
are $f,g$.  More explicitly, put
\[
 \begin{aligned}
 \mathsf B&=(u,Ru,R^2u,R^3u),\\
 d&=(C_f-C_g)e_4=(0,-1,-1,-1)^t,\\
 \mathsf B_c&=(d,C_fd,C_f^2d,C_f^3d).
 \end{aligned}
\]
Here $\det\mathsf B_c=1$.  Uniqueness of the covector in
\eqref{maps:eq:F.9} gives the intertwiner
$\mathsf P=\mathsf B\mathsf B_c^{-1}$, with
$R\mathsf P=\mathsf P C_f$ and $S\mathsf P=\mathsf P C_g$.
Thus $(T_1,T_2)$ is rationally conjugate to
$(C_f^{-1},C_g)$, with the chosen peripheral conventions.

\item Use the lexicographic basis in $\bigwedge^2\mathbb Q^4$.
Its common invariant vector is $(1,1,1,1,1,1)^t$.
Let $Q$ be the identity matrix of size six with first column
replaced by this vector.  It is unimodular.  Direct multiplication
gives the exact identities
\[
 \left(Q^{-1}\bigwedge^2(C_f^{-1})Q\right)_{2:6,2:6}
   =A_{2:6,2:6},\qquad
 \left(Q^{-1}\bigwedge^2(C_g)Q\right)_{2:6,2:6}
   =B_{2:6,2:6}.
 \tag{F.11}\hypertarget{tag-02EH}{}\label{maps:eq:F.11}
\]
They are checked over $\mathbb Z$ in
\texttt{fourfolds\_periods.py}.
The invariant line is the inverse of the polarization, up to a
nonzero rational scalar.  Projection along that line identifies
the quotient with the primitive exterior square $\mathcal W$.

\item The differentials $\omega_1=dx/y$ and $\omega_2=x\,dx/y$ span
the extended Hodge bundle everywhere.  At the node they are
independent dualizing differentials: $\omega_1$ has nonzero
opposite residues on the normalization and
$(x-r)\omega_1$ is a nonzero holomorphic differential there.
The extension assertion can be checked on the nodal annulus.
In coordinates $uv=t$ a relative dualizing differential has
a convergent Laurent expansion
$\sum_n a_n(t)u^n\,du/u$.  Integration on the vanishing circle
is $2\pi i a_0(t)$.  Integration across the annulus has the
single logarithmic term $a_0(t)\log t$; integration of every
nonconstant Laurent monomial and of the complementary fixed
arcs gives a holomorphic remainder.  Thus the period vector
has the form $\exp(sN_H)$ times a holomorphic vector,
$s=\log t/(2\pi i)$.  Untwisting by $\exp(-sN_H)$ is exactly
the nilpotent canonical extension used here.  The two independent
dualizing limits remain a basis of its Hodge subbundle.
Their isotropic pairing, zero away from zero, remains zero
at the limit by holomorphic continuation, so that limit is
Lagrangian.
The root-divisor calculation is \eqref{maps:eq:F.8}.
The nonzero residue of $\omega_1$ also shows
\[
 N_H\ell_\infty\ne0,
 \qquad \ell=\mathcal T^3\subset\mathcal E,
 \tag{F.12}\hypertarget{tag-02EI}{}\label{maps:eq:F.12}
\]
where $\ell_\infty$ denotes its limit at the nodal cusp.
Indeed integration on the vanishing circle tends to
$2\pi i$ times that residue; the transvection sends the cohomology
class of a differential to that period times its vanishing
cohomology vector.

\end{enumerate}
\end{proof}




\begin{theorem}[The integral exterior-square quotient]\hypertarget{tag-02QT}{}\label{night:P1.2}
\StatementRecord{night:P1.2}{proved}{02QT}{the definitions and constructions stated in the proof}
The lower blocks of $A,B$ are precisely the quotient representations of
$\bigwedge^2(C_f^{-1}),\bigwedge^2C_g$ in the basis $Q$.
The integral quotient is the dual of the primitive exterior-square
lattice.  The natural map from the primitive lattice into that quotient
has cokernel $\mathbb Z/2$; these two integral representations are not
isomorphic.  Their rationalizations are isomorphic by that same map.
\end{theorem}
\begin{proof}
\begin{enumerate}[label={$\langle1\rangle$\arabic*.}]
\item The two matrices obtained from their two-by-two minors are
\[
Q^{-1}\bigwedge^2(C_f^{-1})Q=
\begin{pmatrix}
1&-1&0&1&0&0\\0&1&-1&-1&1&0\\0&1&0&-1&0&0\\
0&2&-1&-1&0&1\\0&2&0&-1&0&0\\0&1&1&-1&0&0
\end{pmatrix},
\]
\[
Q^{-1}\bigwedge^2C_gQ=
\begin{pmatrix}
1&0&1&0&0&0\\0&0&-1&0&1&0\\0&0&-1&0&0&1\\
0&0&-1&0&0&0\\0&1&-1&0&0&0\\0&0&-1&1&0&0
\end{pmatrix}.
\]
Their lower five-by-five blocks are the asserted blocks, entry by entry.
Put $L=\bigwedge^2\mathbb Z^4$, let $(\ ,\ )$ be its unimodular wedge
pairing, and set $\omega=(1,1,1,1,1,1)^t$.
Then $(\omega,\omega)=2$ and the primitive sublattice
$P_{\mathbb Z}=\omega^\perp$ is the kernel of
$(1,-1,1,1,-1,1)$.  In the basis with columns
\[
\begin{pmatrix}
1&-1&-1&1&-1\\1&0&0&0&0\\0&1&0&0&0\\
0&0&1&0&0\\0&0&0&1&0\\0&0&0&0&1
\end{pmatrix},
\]
the inclusion followed by projection
$j:P_{\mathbb Z}\longrightarrow L/\mathbb Z\omega$ has matrix
\[
j=\begin{pmatrix}
0&1&1&-1&1\\-1&2&1&-1&1\\-1&1&2&-1&1\\
-1&1&1&0&1\\-1&1&1&-1&2
\end{pmatrix},\qquad \det j=2.
\]
Its Smith factors are $(1,1,1,1,2)$.  The unimodular pairing on $L$
identifies $L/\mathbb Z\omega$ with $P_{\mathbb Z}^\vee$;
under this identification $j$ is the pairing map.

\item We rule out every possible integral intertwiner.  Denote the lower blocks above by $A_5,B_5$.
The Krylov matrix $(e_3,A_5e_3,A_5^2e_3,A_5^3e_3,A_5^4e_3)$
has determinant $-1$, so $A_5$ is cyclic.  Any rational endomorphism
commuting with $A_5$ therefore has the form
$u_0I+u_1A_5+\cdots+u_4A_5^4$.
The entries $(1,2),(1,4),(2,4),(3,1)$ of its commutator with $B_5$
give, on $(u_1,u_2,u_3,u_4)^t$, the coefficient matrix
\[
\begin{pmatrix}
0&1&0&1\\1&-1&0&0\\1&-1&-1&1\\1&1&0&1
\end{pmatrix},\qquad \det=-1.
\]
Thus $u_1=u_2=u_3=u_4=0$, proving that the simultaneous rational
centralizer is exactly $\mathbb QI$.
The natural map $j$ above is equivariant and rationally invertible.
For any rational equivariant map $f:P_{\mathbb Q}\to
(L/\mathbb Z\omega)_{\mathbb Q}$, the composite $fj^{-1}$ lies
in this centralizer.  Consequently $f=cj$ for some $c\in\mathbb Q$.
An integral isomorphism would have determinant $\pm1$, whereas
$\det(cj)=2c^5$.  Its $2$-adic valuation would give
$1+5v_2(c)=0$, impossible because $v_2(c)\in\mathbb Z$.
This proves the asserted failure of integral isomorphism and retains
$j$ itself as the explicit rational isomorphism.

\end{enumerate}
\end{proof}




\begin{theorem}[Integral marking of the pencil]\hypertarget{tag-02QY}{}\label{night:P1.7}
\StatementRecord{night:P1.7}{proved}{02QY}{\SeriesRef{maps:src:F.3.5}, \SeriesRef{night:P1.2}}
The integral $H^1$ monodromy of the pencil has a companion marking with
finite pair $(C_f^{-1},C_g)$, and cusp product $C_fC_g^{-1}$.
Its exterior-square quotient is consequently the integral quotient
representation appearing in \cref{night:P1.2}.
\end{theorem}
\begin{proof}
\begin{enumerate}[label={$\langle1\rangle$\arabic*.}]
\item Let $L=H^1(C_\lambda,\mathbb Z)$ with its integral cup pairing,
and write $T_1,T_2$ for the finite generators of
\cref{maps:src:F.3.5}. Put $A=C_f$ and $B=C_g$.
The rational intertwiner proved there is a map
\[
 q:L_{\mathbb Q}\longrightarrow\mathbb Q^4,
 \qquad qT_1q^{-1}=A^{-1},\quad qT_2q^{-1}=B.
\]
Its relation to the Levelt presentation is
$(h_\infty,h_1,h_0)=(T_1^{-1},T_1T_2,T_2^{-1})$;
the ordered product is one. The known rigidity input and its explicit
rational change of basis are recorded in \cref{maps:src:F.3.5}.
The nodal peripheral is $c=(T_1T_2)^{-1}$. Its primitive
Picard--Lefschetz class $\delta\in L$ and coefficient one are the
geometric node input proved there: the nonseparating vanishing circle
has a transverse integral cycle meeting it once, and the annulus
$xy=t$ turns once. Thus $T_1T_2$ acts by
$v\mapsto v+\epsilon\delta\langle\delta,v\rangle$,
where $\epsilon\in\{1,-1\}$ records the fixed inverse and pairing
conventions.

\item A $B$-invariant alternating form has upper entries
$(u,v,u,u,v,u)$. Invariance under $A$, evaluated on the first and
fourth basis vectors, gives $v=-u$. The common invariant forms are
therefore exactly the multiples of
\[
 J=\begin{pmatrix}
 0&1&-1&1\\-1&0&1&-1\\1&-1&0&1\\-1&1&-1&0
 \end{pmatrix},\qquad
 d=\begin{pmatrix}1\\1\\1\\0\end{pmatrix}.
\]
The Pfaffian of $J$ is $1$, and direct multiplication gives
\[
 A^tJA=B^tJB=J,\qquad
 A^{-1}B-I=d e_4^t,\qquad d^tJ=e_4^t.
\]
Here $e_4$ is the fourth standard vector. Write the transported cup
form as $sJ$ and $q(\delta)=k d$, with $s,k\in\mathbb Q^*$.
Comparison with the Picard--Lefschetz formula gives
$\epsilon s k^2=1$.

\item The integral classes
$\delta_i=T_1^{-i}\delta$, $0\leq i\leq3$, have rational images
equal to the columns of $kD$, where
\[
 D=(d,Ad,A^2d,A^3d)=
 \begin{pmatrix}1&0&-1&0\\1&1&-1&-1\\1&1&0&-1\\0&1&0&0\end{pmatrix},
 \qquad \det D=1.
\]
Their pairing is $\epsilon G$, with
\[
 G=D^tJD=
 \begin{pmatrix}0&1&0&0\\-1&0&1&0\\0&-1&0&1\\0&0&-1&0\end{pmatrix},
 \qquad \det G=1.
\]
These four classes are a basis of $L_{\mathbb Q}$. For any $v\in L$,
write $v=\sum_i x_i\delta_i$ with $x_i\in\mathbb Q$.
The vector $(\langle\delta_i,v\rangle)_i=\epsilon Gx$ is integral.
Since $G^{-1}$ is integral, so is $x$. Thus the four classes are an
integral basis of $L$.

\item The required integral marking is the precise map
$\mu=k^{-1}q:L\longrightarrow\mathbb Z^4$.
Indeed $q(L)=kD\mathbb Z^4=k\mathbb Z^4$. Its inverse is specified
on standard vectors by
\[
 \mu^{-1}(e_j)=\sum_{i=0}^3(D^{-1})_{i+1,j}\delta_i,
 \qquad
 D^{-1}=\begin{pmatrix}1&-1&1&0\\0&0&0&1\\0&-1&1&0\\1&-1&0&1\end{pmatrix}.
\]
This supplies both integral arrows and retains the cup form
$\epsilon J$. In this marking the finite pair is $(A^{-1},B)$.

\item The exact simultaneous companion reversal is
\[
 W=\begin{pmatrix}0&0&0&1\\0&0&1&0\\0&1&0&0\\1&0&0&0\end{pmatrix}.
\]
It satisfies $W^2=I$, $\det W=1$, $WAW=A^{-1}$,
$WBW=B^{-1}$ and $W^tJW=-J$. These identities give the integral comparison of the two companion conventions, including
the pairing sign. For the based cusp used in the statement take
$c'=T_1^{-1}cT_1=(T_2T_1)^{-1}$. Then
\[
 \mu c\mu^{-1}=B^{-1}A,\qquad
 \mu c'\mu^{-1}=A(B^{-1}A)A^{-1}=AB^{-1}.
\]
The homology matrices are the inverse transposes under the evaluation pairing with $L$.

\item Apply $\bigwedge^2\mu$ to the integral wedge lattice. In the
lexicographic basis of \cref{night:P1.2}, the invariant vector is
$\omega=(1,1,1,1,1,1)^t$. Let $Q$ have first column $\omega$ and
remaining columns the last five standard vectors, as in that result.
The induced quotient is the map
\[
 [u\wedge v]\longmapsto
 \operatorname{pr}_{2,\ldots,6}
 Q^{-1}\bigl(\mu(u)\wedge\mu(v)\bigr).
\]
Its kernel before passage to the quotient is precisely the invariant
line. Its inverse is induced by $(\bigwedge^2\mu^{-1})Q$ on the last
five standard vectors. The finite quotient matrices are the lower
five-by-five blocks of
\[
 Q^{-1}(\bigwedge^2 A^{-1})Q,\qquad
 Q^{-1}(\bigwedge^2 B)Q.
\]
They are exactly the integral quotient matrices of
\cref{night:P1.2}, and the cusp follows by the same based word
$(T_2T_1)^{-1}$. The invariant form, Picard--Lefschetz rank-one map,
unimodular orbit matrix, integral inverse, reversal, based peripheral
word and exterior-quotient identities are recomputed exactly in
\href{https://github.com/KokunoYumeto/s6-complex-structure-workbench/blob/83d76022bcd442278b760db9ff438afceed6c6ea/pro_returns/S6_SERIES_GRIND_checkpoint_10_20261120T0211Z/code/current_audits/audit_corrections/code/independent_V_markings_06.py}{\texttt{independent\_V\_markings\_06.py}}.
The integral saturation argument in steps 2--4 supplies the map $\mu$
for the geometric lattice.
\end{enumerate}
\end{proof}




\begin{theorem}[The rank-three bundle $P$]\hypertarget{tag-02QW}{}\label{night:P1.5}
\StatementRecord{night:P1.5}{proved}{02QW}{the definitions and constructions stated in the proof}
The wedge map from
$\ell_0\otimes(\mathcal H_{\mathrm{dR}}/\ell_0)$, followed by
projection modulo the nondegenerate polarization bivector, is
injective and has rank-three image. This image defines the bundle
$P$ used in the period construction.
Moreover
$P\cong\mathcal O\oplus T^2\oplus T^4$,
$H^1(P^\vee)=0$, and $H^0(P^\vee)$ is the one-dimensional
space of contractions. Its inclusion into the rank-five quotient
lifts holomorphically to the rank-six extension.
\end{theorem}
\begin{proof}
On the root curve of orders five and eight, let $\mathcal H_D$ be
the logarithmic de Rham bundle of the pencil. Its Hodge subbundle is
$\mathcal E=T^3\oplus T$, with $\ell=T^3$. The polarization identifies
$\mathcal H_D/\mathcal E$ with $\mathcal E^\vee$. This remains true
at the node: the dualizing differentials have independent limits,
and their isotropic pairing extends there holomorphically.
The primitive projection $p$ and the marked extension give
\begin{align}
 0&\longrightarrow\mathbb C e\longrightarrow\mathcal V
   \longrightarrow\mathcal W\longrightarrow0,
 \hypertarget{tag-02EK}{}\label{maps:eq:F.14}\\
 P&=p(\ell\wedge\mathcal H_D).
 \hypertarget{tag-02EL}{}\label{maps:eq:F.15}
\end{align}
In logarithmic frames the first sequence extends over the cusp,
since the nilpotent exponential commutes with its inclusion and
quotient. It extends equivariantly over the finite root charts.
Every nonzero vector in $\ell\wedge\mathcal H_D$ is decomposable, whereas the polarization bivector has nonzero square. Thus primitive projection is injective there. The Hodge sequence therefore gives
\begin{equation}
 0\longrightarrow T^4\longrightarrow P
   \longrightarrow\mathcal O\oplus T^2\longrightarrow0.
 \hypertarget{tag-02EM}{}\label{maps:eq:F.16}
\end{equation}
For the coarse projection $\rho$, invariant sections in the two
root charts give
\begin{equation}
 \rho_*T^j=\mathcal O_{\mathbb P^1}
 \left(\left\lfloor\frac{2j}{5}\right\rfloor+
       \left\lfloor-\frac{3j}{8}\right\rfloor\right).
 \hypertarget{tag-02EN}{}\label{maps:eq:F.17}
\end{equation}
Finite-group averaging makes descent exact. For $j=\pm2,\pm4$
the degree is $-1$, so both cohomology groups vanish. Consequently
$P=\mathcal O\oplus T^2\oplus T^4$ and
$H^1(P^\vee)=0$. The pullback of \eqref{maps:eq:F.14} to $P$
splits holomorphically. Its splittings differ by $z\alpha e$,
since $H^0(P^\vee)=\mathbb C\alpha$.
\end{proof}

\subsection{The algebraic connection and its splitting}

\begin{theorem}[The algebraic connection]\hypertarget{tag-02QO}{}\label{night:P1.15}
\StatementRecord{night:P1.15}{proved}{02QO}{\SeriesRef{night:P1.1}, \SeriesRef{night:P1.8}}
The rank-six flat extension, as a complex flat bundle with its invariant
line, has the following explicit algebraic connection model.  Set
$\omega_k=x^kdx/y$, $0\leq k\leq3$, and use the quotient frame
\[
(w_2,w_1,w_0,w_{-1},w_{-2})
 =(\omega_{01},\omega_{02},\omega_{03},\omega_{13},\omega_{23}),
\qquad \omega_{12}=-3\omega_{03}.
\]
These are meromorphic frames of
$(L^2,L,\mathcal O,L^{-1},L^{-2})$, where $L=T^2$ on the root
orbifold of orders $5,4$.  If a frame is written as a row of vectors,
its derivative is right multiplication by
\[
\Gamma_6=\begin{pmatrix}0&b\\0&\Gamma_W\end{pmatrix},\qquad
b=\frac1{\lambda-\lambda_*}
 \left(-\frac{25}{6\lambda^2},-\frac{125}{48\lambda},0,1,0\right),
\tag{1.15.1}
\]
where
\[
\resizebox{\textwidth}{!}{$\displaystyle
\Gamma_W=\frac1{3125\lambda+256}
\begin{pmatrix}
-8(625\lambda+48)/\lambda&-8&10\lambda&75\lambda^2/2&0\\
-125/(2\lambda)&-4(625\lambda+48)/\lambda&8&0&75\lambda^2/2\\
300/\lambda^2&-375/\lambda&0&-48&60\lambda\\
120/\lambda^3&0&125/(2\lambda)&4(625\lambda+48)/\lambda&-8\\
0&120/\lambda^3&50/\lambda^2&-125/(2\lambda)&8(625\lambda+48)/\lambda
\end{pmatrix}.$}
\tag{1.15.2}
\]
This model specifies a canonical algebraic splitting, denoted
$\sigma_b$.  The compatible contraction used for the real parameter is
$\alpha=(625\pi i/128)\alpha_{\rm alg}$, where
$\alpha_{\rm alg}(w_0)=4/3$ and it vanishes on $w_1,w_2$.
\end{theorem}
\begin{proof}
\begin{enumerate}[label={$\langle1\rangle$\arabic*.}]
\item Reduction modulo exact differentials is completely explicit.  To
differentiate $\omega_k$, solve the polynomial identity
\[
-\tfrac12x^k\partial_\lambda f
 =f\sum_{j=0}^3G_{kj}x^j+R_k'f-\tfrac12R_kf',
 \qquad\deg R_k\leq4,
\]
where $f=x^5+\lambda^3x+\lambda^4$.  The last two terms are the
numerator of $d(R_k/y)$.  Matching the nine coefficients gives the
four-by-four Gauss--Manin matrix exactly.
At the branch point at infinity, a local coordinate $z=x^{-1/2}$
gives $\omega_k=-2z^{2-2k}(1+O(z^8))\,dz$.
Residues of a primitive times a differential give the alternating
pairings $J_{03}=4/3$, $J_{12}=4$, with their negatives in reverse
order and all other entries zero.  This form is nondegenerate.
The invariant bivector is therefore $3\omega_{03}+\omega_{12}$.
Taking the exterior-square connection and this quotient gives
(1.15.2). The polynomial reduction, all exact primitives, the horizontal
polarization relation and all twenty-five quotient entries are reconstructed
in \href{https://github.com/KokunoYumeto/s6-complex-structure-workbench/blob/83d76022bcd442278b760db9ff438afceed6c6ea/pro_returns/S6_SERIES_GRIND_checkpoint_10_20261120T0211Z/code/current_audits/audit_corrections/code/independent_V_connection_spectra_06.py}{\texttt{independent\_V\_connection\_spectra\_06.py}} and its JSON output.

\item The root-chart powers computed from $x^kdx/y$ are
$u^{6k-9}$ at zero and $u^{6-4k}$ at infinity.
Their nondegenerate pairing shows that the four resulting regular
frames form the full de Rham bundle on those charts.  At the node
the computed connection has a simple pole and square-zero residue;
hence these are the nilpotent logarithmic extension frames there.
Thus the four line bundles are
$T^3,T,T^{-1},T^{-3}$, and their quotient exterior-square lines are
the stated powers of $L$.

\item For the coarse projection $\rho$ on the effective root curve,
\[
 \rho_*L^j=\mathcal O_{\mathbb P^1}
 \left(\left\lfloor\frac{4j}{5}\right\rfloor+
       \left\lfloor-\frac{3j}{4}\right\rfloor\right).
\]
The degree is zero for $j=0$ and minus one for $j=\pm1,\pm2$.
Finite stabilizer invariants are exact, so all five $H^1(L^j)$ vanish.  Consequently every holomorphic
extension of this $W$ by $\mathcal O$ splits as a bundle.
In such a splitting its flat connection is an upper block extension
by a row $b$.  Regularity at the finite root charts and a logarithmic
pole at the cusp require, entry by entry,
\[
b=\frac1{\lambda-\lambda_*}
 (a\lambda^{-2},b'\lambda^{-1},0,c,d\lambda).
\]
For example the $L^j$ entry, after $\lambda=u^4$, has exponent
$3j+3+4p$ for a term $\lambda^p/ (\lambda-\lambda_*)$;
the infinity chart requires $-4j-1-5p\geq0$.
These two inequalities leave exactly the four displayed terms.
The equation that the full cusp residue have square zero is
$b_{\rm res}(\Gamma_W)_{\rm res}=0$.
Its solution space has basis
\[
(-25/6,-125/48,1,0),\qquad(5/3,-25/3,0,1).
\]
The holomorphic splitting changes are the constant multiples of the
row $r_0=(0,0,1,0,0)$ selecting $w_0$.  Its connection coboundary
is the third row $r_0\Gamma_W$, whose parameter vector is
$(-48b_1+60b_2)/3125$.  The quotient by this exact gauge direction
is one-dimensional, with coordinate $s=c+4d/5$.  The first
vector has $s=1$ and rank-three square-zero cusp residue.

The extension is nontrivial, since its dual has no common
invariant vector by \cref{night:P1.1}; consequently its coordinate $s$ is nonzero.
Set $t=3125d/60$.  Its full row is
$\beta=s b+t r_0\Gamma_W$.  The constant bundle maps
\[
 C=\begin{pmatrix}1&-t r_0\\0&sI_5\end{pmatrix},\qquad
 C^{-1}=\begin{pmatrix}1&(t/s)r_0\\0&s^{-1}I_5\end{pmatrix}
\]
are holomorphic in the stated root frames and fix the invariant
section $e$.  Direct multiplication gives
\[
 C^{-1}\Gamma_6 C+C^{-1}dC
   =\begin{pmatrix}0&\beta\\0&\Gamma_W\end{pmatrix}.
\]
This proves the required equivalence of the complete flat bundles
with their fixed invariant line.  Evaluation at the marked basepoint
gives a common intertwiner for both specified meridians.  Thus its
existence follows from the exact connection maps.  

\item For the last normalization, the kernel of the quotient cusp residue
on $(w_2,w_1,w_0)$ is generated by
$v=(-4096/1171875,64/1875,1)^t$.
It satisfies $b_{\rm res}(v)=625/192$ and
$\alpha_{\rm alg}(v)=4/3$.
The logarithmic monodromy is $2\pi i$ times this residue.
We fix the contraction on this kernel line by the exact companion identity $N_V(\widetilde v)=\alpha(v)e$.
Since the space of global contraction forms is one-dimensional, the
ratio $(625/192)/(4/3)=625/256$ fixes its scalar everywhere and gives
the displayed $625\pi i/128$.

\end{enumerate}
\end{proof}

\begin{definition}[The normalized height function]
\hypertarget{tag-PV03}{}\label{def:hb}
\StatementRecord{def:hb}{proved}{PV03}{}
Let $e=e_0$ be the fixed real section of the rank-six system,
and let $\sigma_b:P\to V$ be its algebraic splitting.
The invariant bivector, primitive projection and contraction are
normalized in the companion marking as in \cref{night:P1.15}.
For each regular base point, including the finite root points,
let $n$ be the real vector in $P\cap\bar P$ with $\alpha(n)=1$.
We define $h_b$ by
\begin{equation}
 \sigma_b(n)-\overline{\sigma_b(n)}=2i h_b e.
 \hypertarget{tag-02ER}{}\label{maps:eq:F.20}
\end{equation}
\end{definition}
\begin{proof}[Well-definedness]
Let $u$ be a local nonzero section of $\ell_0$.
Its Hodge metric is positive, so $\operatorname{span}_{\mathbb C}\{u,\bar u\}$
is a nondegenerate symplectic plane. The inverse polarization
bivector $\omega$ has a nonzero component on the exterior square
of its symplectic orthogonal complement. Consequently an equality
$u\wedge v-\bar u\wedge w=t\omega$ forces $t=0$; with $t=0$
its common value belongs to $\mathbb C(u\wedge\bar u)$.
The primitive projection therefore gives
\begin{equation}
 P\cap\bar P=p(\ell_0\wedge\bar\ell_0).
 \hypertarget{tag-02EP}{}\label{maps:eq:F.19}
\end{equation}
We make the scalar in the primitive decomposition explicit.
In the algebraic de Rham frame put
$\omega_{\rm alg}=3\omega_{03}+\omega_{12}$ and let
$\alpha_{\rm alg}(v)$ be the contraction of its unique wedge
lift $j(v)$ with the residue pairing $J_{03}=4/3$, $J_{12}=4$. Then
$\operatorname{contr}_{J}(\omega_{\rm alg})=8$. The raw primitive
coefficient is therefore $\alpha_0=\alpha_{\rm alg}/8$, whereas
\cref{night:P1.15} fixes
\[
 \alpha=\frac{625\pi i}{128}\alpha_{\rm alg}
        =\frac{625\pi i}{16}\alpha_0,
 \qquad
 \omega_{\rm norm}=\frac{16}{625\pi i}\omega_{\rm alg}.
\]
For the primitive representative $v$ of a quotient vector, its
unique wedge lift is exactly
\begin{equation}
 j(v)=v+\alpha_0(v)\omega_{\rm alg}
     =v+\alpha(v)\omega_{\rm norm}.
 \hypertarget{tag-02EO}{}\label{maps:eq:F.18}
\end{equation}

Here $\omega_{\rm norm}$ is real in the marked flat system,
although the algebraic formula contains $i$. To see the precise
comparison, use the companion cup matrix $\varepsilon J_c$ of
\cref{night:P1.7}, with $\varepsilon\in\{1,-1\}$, and its
invariant bivector $\omega_c=(1,1,1,1,1,1)^t$.
Contraction with $J_c$ gives $\operatorname{contr}_{J_c}(\omega_c)=2$.
For the rank-four period comparison $\Pi$, the residue formula is
$\Pi^t(\varepsilon J_c)\Pi=2\pi iJ$.
Indeed, cut the curve along a symplectic polygon and small circles
around the poles; the boundary integral of a primitive times the
second differential is the alternating period pairing on the
polygon and $2\pi i$ times the sum of residues on the circles.
It follows by inversion of the two alternating forms that
$(\bigwedge^2\Pi)\omega_{\rm alg}=(8\pi i/\varepsilon)\omega_c$.

The scaled exterior-square map
$\Psi=(625\varepsilon/256)\bigwedge^2\Pi$ induces the
quotient comparison in the fixed rank-six extension. Its scalar is forced by the cusp contraction identity: on the
kernel vector of \cref{night:P1.15} the extension residue is
$625/192$, while $\alpha_{\rm alg}$ has value $4/3$; the
logarithmic monodromy contributes $2\pi i$. The ratio is
$625\pi i/128$, exactly the pullback of
$\operatorname{contr}_{J_c}$ under this scaled comparison.
The one-dimensional extension calculation in
\cref{night:P1.15} then lifts it uniquely to the normalized
rank-six comparison fixing $e$.
Consequently $\Psi(\omega_{\rm norm})=\omega_c/2$, and
the induced primitive comparison sends $\alpha$ to
$\operatorname{contr}_{J_c}$. Both are real in the companion
marking. For $r=ip(u\wedge\bar u)$, the value $\alpha(r)$
is real and nonzero by the positive polarization pairing.
Thus $n=r/\alpha(r)$ is the unique real vector specified above.
Multiplying $u$ by a nonzero complex number multiplies both
$r$ and $\alpha(r)$ by its squared modulus, and leaves $n$
unchanged.

The image of $\sigma_b(n)$ in the quotient is real. Its
difference from its conjugate is therefore an imaginary
multiple of the real vector $e$, which gives a unique real
$h_b$. Every marked monodromy is real and fixes $e$, so the
same equality defines one function on the base. The construction
is real analytic on each root chart; it is invariant under its
finite deck action and hence descends.

For its coordinate expression, let $H'=H\Gamma_6$, $H(1)=I$.
Take the positive zero meridian with stem from $1$ to $1/32$
and a counterclockwise circle of radius $1/32$, and the infinity
meridian $\lambda=\exp(-2\pi it)$, $0\le t\le1$.
Let $H_0,H_\infty^{\rm mon}$ be their based continuations.
The comparison matrix $X$ is specified by
\[
 XH_0=BX,\qquad XH_\infty^{\rm mon}=AX,\qquad Xe_0=e_0.
\]
Its existence is the marked flat-bundle identification in
\cref{night:P1.15}. The equations $RA=AR$, $RB=BR$ for the
explicit integral matrices force $R$ to be scalar. Thus any
two such comparisons differ by a scalar, and $Xe_0=e_0$ fixes
that scalar. This proves uniqueness with the based meridians
specified.

Put $Y=XH$ and $k=625\pi/96$. Since
$\alpha(w_0)=ki$ and $\alpha(w_1)=\alpha(w_2)=0$,
the vector $kn$ has coordinates
$(0,v_2,v_1,-i,0,0)^t$ in the quotient lift for a unique
pair $(v_2,v_1)$. Equivalently, the lower five coordinates
of its image under $Y$ are real. Equation
\eqref{maps:eq:F.20} gives
\[
 h_b(\lambda)=\frac{96}{625\pi}
 \operatorname{Im}\bigl(Y_{01}v_2+Y_{02}v_1-iY_{03}\bigr).
\]
This proves the existence, uniqueness and normalization of the
real system used in the height computation.
\end{proof}

\begin{remark}[Change of splitting]
\hypertarget{tag-PV04}{}\label{rem:V:splitting-change}
Every holomorphic splitting of the pullback extension over $P$
has the form $\sigma_b-a\alpha e$, $a\in\mathbb C$, because
$H^0(P^\vee)=\mathbb C\alpha$. Its parameter and height are
related to the ones fixed here by the exact equalities
\[
 (\sigma_b-a\alpha e+\xi_{\rm alt}\alpha e)(P)
  =(\sigma_b+(\xi_{\rm alt}-a)\alpha e)(P),
 \qquad h_{\rm alt}=h_b-\operatorname{Im}a.
\]
Thus $\xi=\xi_{\rm alt}-a$. We use $\sigma_b$ and $\xi$
throughout.
\end{remark}

\subsection{The lifted periods and compactification}

\begin{theorem}[Lifted periods and their exact parameter domain]
\hypertarget{tag-02O8}{}\hypertarget{tag-02QZ}{}
\label{maps:src:F.3.6}\label{night:P1.8}
\StatementRecord{night:P1.8}{proved}{02QZ}{\SeriesRef{night:P1.5}, \SeriesRef{night:P1.15}}
Let $P$ be the rank-three wedge bundle of \cref{night:P1.5}, and use
its normalized contraction $\alpha$ and the algebraic splitting $\sigma_b$.
The planes
\begin{equation}
 F_\xi=(\sigma_b+\xi\alpha e)(P),\qquad \xi\in\mathbb C,
 \hypertarget{tag-02ES}{}\label{maps:eq:F.21}
\end{equation}
are holomorphic rank-three subbundles of the marked rank-six system.
They extend over its finite root charts of orders five and four.
The function $h_b$ of \cref{def:hb} has an attained finite minimum
$h_{b,\min}$, tends to $+\infty$ at the cusp, and satisfies
\begin{equation}
 F_\xi\cap\overline{F_\xi}=0
 \quad\Longleftrightarrow\quad
 h_b+\operatorname{Im}\xi\ne0
 \hypertarget{tag-02ET}{}\label{maps:eq:F.22}
\end{equation}
at each nonsingular base point, including the finite root points.
Thus the parameters transverse on the completed base are exactly
$\operatorname{Im}\xi>-h_{b,\min}$.
In an integral marking at the cusp their periods are
\begin{equation}
 Z(q)=\frac{\log q}{2\pi i}I_3+C_\xi(q),
 \qquad C_\xi\text{ holomorphic at }q=0.
 \hypertarget{tag-02EJ}{}\label{maps:eq:F.13}
\end{equation}
\end{theorem}
\begin{proof}
The holomorphic splitting and the one-dimensional space of its changes are given by \cref{night:P1.5}. We use the algebraic choice $\sigma_b$ of \cref{night:P1.15}.

The real-intersection calculation and well-definedness proof of
\cref{def:hb} give $P\cap\bar P=\mathbb Cn$ with
$\alpha(n)=1$ and
$\sigma_b(n)-\overline{\sigma_b(n)}=2ih_b e$.
The lifted vectors for parameter $\xi$ differ from their
conjugates by $2i(h_b+\operatorname{Im}\xi)e$.
An arbitrary complex multiple of $n$ gives no further
intersection direction. This proves \eqref{maps:eq:F.22}.

We calculate its cusp behavior in the fixed companion marking.
For a local nonzero $u\in\ell_0$, put $r=ip(u\wedge\bar u)$
and $a_r=\alpha(r)$, so that $n=r/a_r$.
Put $D=C_fC_g^{-1}-I$, $N_V=(BA)^{-1}-I$ and let $N_W$ be its
quotient. The only nonzero column of $D$ is
$(0,1,1,1)^t$, and the top row of $N_V$ is
$(0,0,0,-1,1,-1)$. The wedge projection and normalized contraction
satisfy, for all $a,b\in\mathbb C^4$,
\[
 (N_V)_{0,*}(0,p(a\wedge Db))
 =\operatorname{contr}(a\wedge Db)=a_0b_0,
 \qquad N_Wp(a\wedge Db)=0.
\]
These identities follow by expanding the two-by-two minors in
the companion basis. Write $u=\exp(sD)a(q)$ with
$s=\log q/(2\pi i)$ and $y=\operatorname{Im}s$.
The vector $a$ is holomorphic, and $a_0(0)\ne0$ because
$dx/y$ has nonzero nodal residue. In the logarithmic frame,
\[
 e^{-sN_W}r=ip(a\wedge\bar a)+2y\,p(a\wedge D\bar a),
 \qquad
 a_r=i\operatorname{contr}(a\wedge\bar a)+2y|a_0|^2.
\]
The splitting is bounded in that frame. On returning by
$e^{sN_V}$, its imaginary invariant-coordinate contribution is
$2|a_0|^2y^2+O(y)$. Division by
$a_r=2|a_0|^2y+O(1)$ yields
$h_b(q)=y+O(1)$, uniformly with $\operatorname{Re}s$ in a unit
interval. Thus $h_b(q)\to+\infty$ at the cusp.

The same identities give the logarithmic graph. At the limiting
plane choose a symplectic basis $a,b,d_1,d_2$ with $N_Hb=a$
and $\ell=\mathbb Cb$. Then
\begin{equation}
 P_0=\langle p(b\wedge a),b\wedge d_1,b\wedge d_2\rangle,
 \quad N_W(b\wedge d_i)=a\wedge d_i,
 \quad N_Wp(b\wedge a)=0.
 \hypertarget{tag-02EU}{}\label{maps:eq:F.23}
\end{equation}
The first generator maps under $N_V$ to the nonzero invariant
line, while the other two map to independent quotient directions.
Therefore $N_V:F_{\xi,0}\to\ker N_V$ is an isomorphism for
every $\xi$. In the primitive integral Jordan marking, this
makes the untwisted plane a graph with holomorphic coefficient
$C_\xi(q)$. Restoring $e^{sN_V}$ adds $sI_3$, giving
\eqref{maps:eq:F.13}.

The height extends continuously over both finite root charts.
Its cusp growth and compactness away from that cusp give an
attained minimum. The completed punctured base is connected, so
its image is $[h_{b,\min},\infty)$. Equation
\eqref{maps:eq:F.22} gives the exact parameter half-plane.
The fourth power of the order-eight action on $\mathcal H_D$
is $-I$; it acts trivially on $\mathcal W,P,\alpha$ and
$\mathcal V$. The equivariant lifted plane therefore descends
to the order-four chart. The other chart has order five.
\end{proof}












\begin{theorem}[compact holomorphic filling]\hypertarget{tag-02O9}{}\label{maps:src:F.3.7}\hypertarget{tag-02QJ}{}\label{night:P1.10}
\StatementRecord{maps:src:F.3.7}{proved}{02O9}{\SeriesRef{maps:src:F.2.6}, \SeriesRef{maps:src:F.2.5}, \eqref{maps:eq:F.24}, \eqref{maps:eq:F.13}, \eqref{maps:eq:E.4}, \eqref{maps:eq:E.5}, \SeriesRef{maps:src:E.1.2}, \SeriesRef{maps:src:E.1.3}}

For every transverse period parameter of \cref{night:P1.8}, the
marked gluing of \cref{pub:V:smooth-gluing} has a compact smooth
complex realization and a proper holomorphic map to $\mathbb P^1$.
Its regular fibres are complex three-tori; its two finite reduced
fibres are the free affine quotients of orders five and four in
\cref{pub:V:rank-six-data}; its remaining fibre is the index-one
Kuhn degeneration. With its complex orientation its fundamental
group is trivial and its Euler number is six.

\end{theorem}


\begin{proof}
\begin{enumerate}[label={$\langle1\rangle$\arabic*.}]
\item We first fix the dual convention.  The preceding construction
has cohomology monodromies $A,B$; \cref{pub:V:rank-six-data} uses homology
monodromies $A^t,B^t$.  All three special values are real.
Choose a positive real basepoint and paths along the positive
real circle to infinity and to zero.  Complex conjugation fixes
these paths and reverses the respective peripheral meridians.
Pull back the preceding period family by conjugation of the
base and conjugate its complex fibres.  Its Hodge subbundle is
holomorphic, with local expression
$\overline{\mathcal F(\bar\lambda)}$, and its cohomology
monodromies are $A^{-1},B^{-1}$.  Its homology monodromies are
therefore exactly $A^t,B^t$.
Put $M_1=A^t$, $M_2=B^t$, $M_o=M_2M_1$, $M_0=(M_1M_2)^{-1}$ and
$J=\operatorname{diag}(I_3,-I_3)$.  The based peripheral word is
$\kappa(g_0)=g_1^{-1}g_0^{-1}g_1$, so
$M_0=M_1M_o^{-1}M_1^{-1}$.  For the target marking $S$,
direct integral multiplication gives
\[
 S^{-1}M_0S=\begin{pmatrix}I_3&I_3\\0&I_3\end{pmatrix},\qquad
 S_o=M_1^{-1}SJ,\qquad
 S_o^{-1}M_oS_o=\begin{pmatrix}I_3&I_3\\0&I_3\end{pmatrix}.
\]
Any two markings of this positive standard form differ by
$G=\left(\begin{smallmatrix}P&Q\\0&P\end{smallmatrix}\right)$,
where $P\in\operatorname{GL}_3(\mathbb Z)$ and $Q$ is integral:
this follows by multiplying $G$ with the displayed unipotent
matrix.  The period change is
\[
 (I_3,Z)G=(P,Q+ZP),\qquad
 P^{-1}(P,Q+ZP)
  =(I_3,sI_3+P^{-1}Q+P^{-1}CP).
\]
Thus the bounded holomorphic logarithmic form holds in $S_o$.
In the real-normalized cusp coordinate,
$q_o=\bar q$ and $s_o=-\bar s$.
Base and fibre conjugation therefore give
$(I_3,-sI_3+C^\kappa(q))$, with
$C^\kappa(q)=\overline{C(\bar q)}$.
The period-cycle change $J$ gives the target marking
$M_1S_oJ=S$ and the positive periods
$(I_3,sI_3-C^\kappa(q))$.


\item For a holomorphic Hodge subbundle $\mathcal F$ with
$\mathcal F\cap\bar{\mathcal F}=0$, restriction of integral
homology functionals embeds the homology lattice discretely in
$\mathcal F^*$.  Real transversality makes its real span the
whole real vector space.  The quotient
$\mathcal F^*/\mathcal V_{\mathbb Z}^*$ is consequently a
holomorphic family of compact complex three-tori over $U$,
with its zero section.  This also proves local properness,
by a smoothly varying fundamental parallelepiped.

\item On the root disc at the finite point of order $m_i$, the family
extends smoothly and its deck map has linear part $M_i$.
The invariant vector $v_i$ of \cref{pub:V:rank-six-data} gives a holomorphic
torsion section $v_i/m_i$.  Use the diagonal affine deck action
\[
 (z,s)\longmapsto(M_i z+v_i/m_i,\zeta_{m_i}s).
 \tag{F.24}\hypertarget{tag-02EV}{}\label{maps:eq:F.24}
\]
Here $M_i z$ denotes the induced complex-linear map
$\mathcal F_s^*\to\mathcal F_{\zeta_{m_i}s}^*$ from the marked
lattice action, not multiplication of a complex three-vector
by a six-by-six matrix.  Likewise $v_i/m_i$ means its evaluation
as a point of the target torus.
It has order $m_i$, since the $m_i$th translation is integral.
Every nonidentity power is fixed-point-free on the central
torus: the invariant integral covector $\ell$ takes $v_i$ to
$1$ or $-1$, so a fixed point would require the noninteger
$\pm j/m_i$ to vanish in $\mathbb R/\mathbb Z$.
Away from the central fibre the base rotation is already free.
The quotient is therefore smooth and proper over the coarse
disc, with precisely the required multiple fibre.

\item Its marked punctured family is holomorphically identified with
the original torus family by logarithmic translation.  On the
root cover use $r(s)=(\log s/(2\pi i))v_i$.  Then
$r(\zeta_{m_i}s)-M_i r(s)=v_i/m_i$, modulo the integral lattice;
the change of branch of $\log s$ is itself an integral vector.
Thus translation by $r$ conjugates the linear deck action to
\eqref{maps:eq:F.24}.  This fixes the affine clutching, including its sign.

\item At the cusp use the integral marking of \eqref{maps:eq:F.13}, the fan \eqref{maps:eq:E.4}
and the action \eqref{maps:eq:E.5} with $C_3=C$.
Its maximal charts are $\mathbb C^4$ with map
$t=z_0z_1z_2z_3$, and its transition maps are the toric changes
of coordinates multiplied by the holomorphic units
$\exp(2\pi i C(t)\lambda)$.
The multiplicative comparison is
\[
 q=\bar q_o,\qquad x=(\bar x_o)^{-1},\qquad
 \nu=-\nu_o,
 \qquad C_{\rm new}(q)=-C^\kappa(q).
\]
It carries the old period action to
$x\mapsto\exp(2\pi iC_{\rm new}(q)\nu)q^\nu x$.
On cocharacters it is
$\Phi=\operatorname{diag}(-I_3,1)$ together with conjugation.
The Kuhn simplex at $v$ with ordering $\sigma$ is carried to
the simplex at $-v-(1,1,1)$ with reversed ordering.
If $N_o$ and $N$ are their height-one generator matrices and
$R$ reverses four columns, then
$N=\Phi N_oR$ and $N^{-1}\Phi=RN_o^{-1}$.
Consequently the full affine-chart map is
$z_j=\overline{z_{o,3-j}}$.  These character identities also
give equality on every monomial overlap; the product of the
four coordinates is the conjugate base coordinate.
The already displayed period-action equality makes the map
descend to the analytic period quotient.  Hence the fan and
marking in this comparison are precisely the prescribed ones.

The local filling theorem \cref{night:P1.9} applies to this bounded
holomorphic correction. It gives smoothness, proper discontinuity
and properness in the prescribed marking. We glue with zero cusp
translation.

\item The three punctured overlap maps are holomorphic isomorphisms
with the same marked torus family, so their cocycle condition
is literal.  Proper local models and a compact complementary
base give a Hausdorff compact total space.  Every chart is
smooth.  The affine vectors, cusp fan and clutching are exactly
those in \cref{pub:V:smooth-gluing}.  Varying their bounded corrections
and varying the fibre complex structures along a smooth path
does not change these marked smooth fillings: on the finite
root models the underlying real torus family is the fixed
local-system bundle; on the cusp use the marked, stratum-preserving comparison of
\cref{pub:V:cusp-smooth-dependence}.  Thus the resulting underlying smooth manifold
is the gluing of \cref{pub:V:smooth-gluing}. Its fundamental group is trivial and its Euler number is six by that lemma.

\end{enumerate}
\end{proof}

\begin{remark}[The parameter under base and fibre conjugation]
\hypertarget{tag-PV05}{}\label{rem:V:conjugation-parameter}
Put $k=625\pi/96$. In the normalized integral frame the lifted
period columns are
$W_\xi=(Y_{*,1},Y_{*,2},Y_{*,3}+ki\xi e_0)$.
The rational connection and the based meridians give the real
matrix $K=X\bar X^{-1}$ with
$K^2=I$, $Ke_0=e_0$ and
$Y(\bar\lambda)=K\overline{Y(\lambda)}$.
Consequently the simultaneous base and fibre conjugation has
exact parameter map $\xi\mapsto-\bar\xi$:
\[
 \overline{W_{-\bar\xi}(\bar\lambda)}=K W_\xi(\lambda).
\]
This map preserves $\operatorname{Im}\xi>c_*$ and fixes every
$\xi=ic$, $c\in\mathbb R$. Applying both conjugations with this
parameter map gives a holomorphic family in the coordinate
$\xi$, with cohomology matrices $A^{-1},B^{-1}$ and homology
matrices $A^t,B^t$. Its finite and cusp fillings are the ones
specified in \cref{maps:src:F.3.7}.
\end{remark}




\section{The integral eight-manifold}






  






\begin{theorem}[Homology of the finite fillings]\hypertarget{tag-02QU}{}\label{night:P1.3}
\StatementRecord{night:P1.3}{proved}{02QU}{the definitions and constructions stated in the proof}
The two finite filling pieces have first integral homology
$\mathbb Z^2$ and $\mathbb Z^2\oplus\mathbb Z/2$.
\end{theorem}
\begin{proof}
\begin{enumerate}[label={$\langle1\rangle$\arabic*.}]
\item For $K_i=T^6/\langle\tau_i\rangle$, the universal covering deck group
has commuting translations $x^u$, $u\in\Lambda$, and a generator $a_i$,
with the complete relations
\[
a_ix^ua_i^{-1}=x^{M_iu},\qquad a_i^{m_i}=x^{v_i}.
\]
Indeed the affine lifts satisfy these relations, and they reduce every
word uniquely to $x^ua_i^r$, $0\leq r<m_i$.
The abelianization relation matrix has columns
$((M_i-I)e_j,0)$ and $(-v_i,m_i)$.
Smith reduction of these explicit seven-by-seven matrices gives the
groups above.  The disk quotient
$(T^6\times\Delta)/\langle(\tau_i,\zeta_{m_i})\rangle$
retracts onto $K_i$ by equivariant radial contraction.  Freeness supplies
ordinary smooth quotient charts by disjoint translates of small balls.

\end{enumerate}
\end{proof}


\begin{lemma}[The integral restriction complex]
\hypertarget{tag-PV12}{}\label{pub:V:integral-restrictions}
\StatementRecord{pub:V:integral-restrictions}{proved}{PV12}{}
For the marked smooth gluing of \cref{pub:V:smooth-gluing}, the
integral image sheaves $\mathcal S_q$ and the point-supported
restriction kernels have the groups computed below. For $1\le q\le4$,
$H^2(\mathcal S_q)=\mathbb Z/2$; the only nonzero groups in base
degrees zero and one are $H^0(\mathcal S_1)=\mathbb Z$ and
$H^1(\mathcal S_3)=\mathbb Z^3$. Each of the three order-four
torsion suspension generators has nonzero image in the corresponding
global coinvariant group.
\end{lemma}
\begin{proof}
\begin{enumerate}[label=\arabic*.,leftmargin=*]
\item We use the marked construction of \cref{maps:src:F.2.5} and the
integral cusp-collapse assertion of \cref{night:P1.11}.
Explicitly, $M_1=A^t$, $M_2=B^t$ and the vectors $v_1,v_2$ are
those printed in item~1 of the proof of \cref{maps:src:F.2.5};
$S$ is the displayed integral marking in item~1 of the proof of
\cref{night:P1.1}.  Write
$m_1=5$, $m_2=4$, and put
\[
\begin{split}
C_1&=\begin{pmatrix}
-1&-1&0&-1&0\\0&1&0&1&1\\0&1&0&0&0\\
0&-1&-1&0&-1\\0&0&1&0&0
\end{pmatrix},\\[2pt]
C_2&=\begin{pmatrix}
-1&0&-1&0&-1\\0&1&0&0&0\\0&0&0&1&0\\
0&-1&0&-1&1\\0&0&1&0&0
\end{pmatrix}.
\end{split}
\]
For the translation vectors $v_i$, the matrices
$P_i=[v_i,(0;C_i)]$ are unimodular.  Direct multiplication gives
\[
 P_i^{-1}M_iP_i=1\oplus D_i,\qquad Q_i=P_i^{-t},
\]
where $D_1$ sends $e_j$ to $e_{j+1\bmod5}$, and $D_2$ sends
$e_j$ to $e_{j+1\bmod4}$ for $0\le j<4$ and $e_4$ to $-e_4$.
These equalities give the integral coordinate maps.  In these
coordinates the affine action is
\[
 (u,y,z)\longmapsto
 (u+1/m_i,D_i y,e^{2\pi i/m_i}z).
\]
Thus its core $F_i$ is the mapping torus of $D_i:T^5\to T^5$.
The function $W=e^{-2\pi i u}z$ trivializes its normal disc bundle.
With $\theta=m_i u\bmod1$, the coarse-base function on that bundle is
\[
                  q=W^{m_i}e^{2\pi i\theta}.
\]
In particular the boundary is $F_i\times S^1_W$ with this specified
map to the base circle.

\item We compute the integral restriction images on this model.
Let $V=\mathbb Z^5$ be the vertical cohomology lattice, let
$T_{i,k}=\bigwedge^kD_i^t$, and let
$N_{i,k}=\sum_{a=0}^{m_i-1}T_{i,k}^{a}$.
The signed permutations preserve the product cell structure with one
vertex and one edge on each circle.  Its cellular differential is zero.
Fibre cells and their prisms give the mapping-torus complex
\[
 C^q(F_i)=\bigwedge^qV\oplus\bigwedge^{q-1}V,
 \qquad d(a,b)=(0,(T_{i,q}-I)a).
\]
The $m_i$-fold covering by the generic six-torus pulls back
$(a,b)$ to $(a,N_{i,q-1}b)$: the horizontal cell is the union of the
$m_i$ consecutive prisms, with their transported coefficients.
Consequently its image in $L_q=\bigwedge^q\mathbb Z^6$ is
\[
 J_{i,q}=\bigwedge^q Q_i
 \left(\operatorname{Inv}T_{i,q}\ \oplus\
 e_0\wedge N_{i,q-1}\bigwedge^{q-1}V\right).
\]
Here vertical coordinates occupy slots $1,\ldots,5$ in the six-lattice.
This formula specifies both the lattice and its inclusion.

We fix its column basis as follows.  Exterior coordinates are in
increasing lexicographic order, with indices starting at zero.
For each signed-permutation orbit, start at its least coordinate
and sum its successive signed images until it returns there.
A positive orbit of length $l$ gives a primitive invariant column;
the norm image has the column $m_i/l$ times that column.
A negative orbit has zero invariant lattice and one $\mathbb Z/2$
in its coinvariant group.  Indeed the successive relations identify
all its coordinates with signed copies of the first, and the last
relation is twice that first coordinate.  Order the positive orbits
by their least coordinates, with vertical columns before horizontal
columns.  For $D_1$ there is no torsion in vertical degrees $1$ through
$4$.  For $D_2$ there is one negative orbit in each such degree; the
positive-orbit counts are $1,2,2,1$.

\item At the cusp, \cref{night:P1.11} supplies a primitive collapse
image and free homology, with the only collapse cokernel equal to
$\mathbb Z$ in degree four.  Dualizing those integral maps therefore
gives
\[
 J_{0,q}=\operatorname{Inv}\bigl(\bigwedge^q M_0^t\bigr),
\]
with cohomological restriction kernel zero for $q=1,2,3$ and
$\mathbb Z$ for $q=4$.
Here are explicit primitive bases, which also fix the column order
in the calculation below.  In the cusp marking $S$, write
the cohomology basis as $x_0,x_1,x_2,y_0,y_1,y_2$; its action sends
$x_i$ to $x_i+y_i$ and fixes $y_i$.
In degree one take $(y_0,y_1,y_2)$.  In degree two take, in order,
\[
 y_i\wedge y_j\ (i<j),\qquad
 x_i\wedge y_i\ (0\le i\le2),\qquad
 x_i\wedge y_j+x_j\wedge y_i\ (i<j).
\]
In degree three take $y_0\wedge y_1\wedge y_2$, then
$x_i\wedge y_j\wedge y_k$ in lexicographic $(j,k)$ order with
$j<k$ and $i=j,k$, and then
\[
\begin{split}
 &x_0\wedge y_1\wedge y_2+x_1\wedge y_0\wedge y_2,\\
 &x_0\wedge y_1\wedge y_2-x_2\wedge y_0\wedge y_1.
\end{split}
\]
Their coordinate equations show that these bases are primitive and
have ranks $3,9,9$.  For degree four set
$J=\left(\begin{smallmatrix}0&I_3\\-I_3&0\end{smallmatrix}\right)$,
and let $H(e_I)$ be the complementary exterior basis vector, with
the sign of the permutation $(I,I^c)$.  Apply
$H\bigwedge^2J$ to the displayed degree-two basis.
The identity $JU=U^{-t}J$ gives the invariant equation, and the
signed-permutation maps $H$ and $\bigwedge^2J$ preserve primitivity.
Transport all these columns by $\bigwedge^qS^{-t}$ to obtain
$J_{0,q}$ in the original marking.  Thus its degree-four rank is nine.

\item Let $\mathcal R_q=R^qf_*\mathbb Z$ on the coarse base $S^2$,
let $j$ include its thrice-punctured locus, and let $\mathcal S_q$
be the image of $\mathcal R_q\to j_*L_q$.
Since $M$ is compact and the base is Hausdorff, the proper-base-change
map identifies its stalks with fibre cohomology
\cite[Tag 09V5, Lemma 20.18.1; chapter PDF p.~30]{StacksProject}.
Its three stalk images are exactly the lattices above.  The kernel
$\mathcal K_q$ is supported at those points: for $q=2,3$ it is the
order-four $\mathbb Z/2$, and for $q=4$ it is that $\mathbb Z/2$
together with the cusp $\mathbb Z$.
Point-supported sheaves have no positive cohomology, so
$H^p(\mathcal R_q)=H^p(\mathcal S_q)$ for $p>0$.

Put $A_q=\bigwedge^qM_1^t$ and $B_q=\bigwedge^qM_2^t$.
The based marking gives
$\bigwedge^qM_0^t=(B_qA_q)^{-1}$.
Oriented cellular duality on the punctured sphere and its two-loop
spine gives the compact-support complex $L_q^2\to L_q$ in degrees
one and two.  Attaching the three stalk lattices gives the full complex
\[
\begin{gathered}
 J_{1,q}\oplus J_{2,q}\oplus J_{0,q}
 \xrightarrow{d_0}L_q^2\xrightarrow{\Delta_q}L_q,\\
 d_0(j_1,j_2,j_0)=(j_1-j_0,j_2-A_qj_0),\\
 \Delta_q(a,b)=(A_q-I)a+(B_q-I)b.
\end{gathered}
\]
For the cusp column, the two consecutive peripheral paths carry
$j_0$ and $A_qj_0$; their boundary is
$(B_qA_q-I)j_0=0$.  The reverse cusp orientation gives their negatives.
The two finite columns are the corresponding single peripheral
loops.  This is the relative cellular sequence for extension by zero
and the three stalk inclusions, so it retains their integral indices.
It also verifies $\Delta_qd_0=0$ with the displayed transports.

Here are finite determinant certificates for its cohomology.  They
use precisely the ordered bases just given, so every matrix is
specified by the marked $M_i,S,C_i$ and exterior minors.
For each $q$ select the following columns $\mathcal P_q$ of
$\Delta_q$ and call their square matrix $E_q$:
\[
\begin{aligned}
\mathcal P_1&=(1,2,5,8,9,10),\\
\mathcal P_2&=(0,1,2,8,9,11,12,13,16,17,23,24,25,26,28),\\
\mathcal P_3&=(1,2,3,5,6,10,11,12,13,15,16,18,19,\\
             &\hspace{22mm}26,29,31,32,34,35,37),\\
\mathcal P_4&=(0,1,5,7,9,10,11,13,16,17,20,23,26,28,29).
\end{aligned}
\]
Direct integer multiplication gives
\[
 \det E_q=(2,2,2,-2)_q,\qquad
 E_q^{-1}\Delta_q\in\operatorname{Mat}(\mathbb Z).
\]
Hence $E_q\mathbb Z^{\binom6q}$ is the entire image of
$\Delta_q$ and its cokernel is $\mathbb Z/2$.
Let $\mathcal F_q$ be the increasing complement of $\mathcal P_q$
in $\{0,\ldots,2\binom6q-1\}$.
These free-coordinate rows identify $\ker\Delta_q$ integrally:
the pivot coordinates of a kernel vector are
$-E_q^{-1}(\Delta_q)_{\mathcal F_q}$ times its free coordinates.
The relation matrix is therefore $R_q=(d_0)_{\mathcal F_q}$.
Its sizes, ranks and primitive minors are
\[
\begin{array}{c|c|c|l}
q&\text{size of }R_q&\operatorname{rank}R_q&
\text{maximal minors}\\ \hline
1&6\times7&6&-20,\ -5,\ 4\\
2&15\times15&15&\det R_2=1\\
3&20\times17&17&-1\\
4&15\times15&15&\det R_4=-1
\end{array}
\]
For $q=1$ the three column sets are respectively
$(0,1,2,3,5,6)$, $(0,1,2,4,5,6)$, $(0,2,3,4,5,6)$,
using all rows.  For $q=3$ use all columns and rows
$(0,1,\ldots,13,15,16,17)$.
These determinants have greatest common divisor one, so the
quotient by the relations is free.  The stalk column counts are
$(2,2,3)$, $(3,3,9)$, $(4,4,9)$ and $(3,3,9)$.
Thus the kernel dimensions and the quotient ranks give exactly
\[
\begin{array}{c|ccc}
q&H^0(\mathcal S_q)&H^1(\mathcal S_q)&H^2(\mathcal S_q)\\ \hline
1&\mathbb Z&0&\mathbb Z/2\\
2&0&0&\mathbb Z/2\\
3&0&\mathbb Z^3&\mathbb Z/2\\
4&0&0&\mathbb Z/2.
\end{array}
\]

\item For $k=1,2,3$ the negative signed orbits have representatives
$e_4$, $e_0\wedge e_2$, $e_0\wedge e_2\wedge e_4$.
Map them by the vertical inclusion followed by $\bigwedge^kQ_2$;
call the resulting columns $w_k\in L_k$.
The following rows $\rho_k$, read modulo two, annihilate
$\Delta_k$:
\[
\begin{aligned}
\rho_1&=(0,1,1,1,1,1),\\
\rho_2&=(1,1,1,1,1,1,1,0,0,1,0,1,0,1,1),\\
\rho_3&=(0,0,0,0,0,0,0,0,0,0,1,1,0,1,0,1,1,1,0,0).
\end{aligned}
\]
The explicit representatives and integer witnesses are
\[
\begin{aligned}
w_1&=(0,0,1,1,1,0)^t,\\
w_2&=(0,0,0,0,0,0,0,0,-1,0,0,-1,0,-1,-1)^t,\\
w_3&=(0,0,0,0,0,0,0,0,0,0,0,0,1,0,1,1,0,0,0,0)^t,\\
z_1&=(0,0,0,-1,-1,-1)^t,\\
z_2&=(-1,0,1,-2,-2,-2,4,2,-1,-3,2,0,-2,0,2)^t,\\
z_3&=(-2,-2,-3,0,0,-1,-1,-1,2,2,-1,3,-1,4,3,\\
   &\hspace{23mm}-3,-2,1,-4,1)^t.
\end{aligned}
\]
They satisfy $E_kz_k=2w_k$ and $\rho_kw_k=1\bmod2$.
The nonzero local torsion class therefore maps to the nonzero global
class in each degree $q=2,3,4$.

\end{enumerate}
\end{proof}

\begin{lemma}[The torsion transgression at the order-four filling]
\hypertarget{tag-02QF}{}\label{mapsaudit:order-four-torsion-transgression}
\StatementRecord{mapsaudit:order-four-torsion-transgression}{proved}{02QF}{\SeriesRef{maps:src:F.2.5}, \SeriesRef{maps:src:F.2.6}, \eqref{maps:eq:F.6}}

Use the order-four filling and the marked matrices of \cref{pub:V:rank-six-data}.
Let $T=\bigwedge^{q-1}D^t$ on the integral vertical cochains of the
five-torus. If $w$ represents a torsion element of
$\operatorname{coker}(T-I)$, its point-supported degree-$q$ class has
Leray differential equal, up to the one boundary-orientation sign, to
the image of $w$ in the global degree-$(q-1)$ coinvariants.  For each of $q=2,3,4$, the unique nonzero local torsion class maps
to the nonzero class in the global $\mathbb Z/2$ quotient.
\end{lemma}


\begin{proof}
\begin{enumerate}[label={$\langle1\rangle$\arabic*.}]
\item Use the explicit matrix $C_2$ in the proof of
\cref{pub:V:integral-restrictions}.  In the original vertical marking $D$ is the
bottom-right block of $M_2$; multiplication gives
$DC_2=C_2D_2$, where $D_2$ is the four-cycle followed by the sign
on its fifth coordinate.  The change of cochain coordinates is
$\bigwedge^kC_2^t$, and the inverse is $\bigwedge^kC_2^{-t}$.
Thus the product cell structure of the signed-permutation torus,
transported by this integral torus isomorphism, has zero cellular
differential and the stated action $T=\bigwedge^{q-1}D^t$.
The fibre cells and their prisms give the mapping-torus
cochain complex.  The degree-$q$ suspension classes are consequently
the coinvariants of $T$ on vertical degree-$q-1$ cochains.

\item Write $u=-x_0$ in the original marking.  The affine action is
$(u,y,z)\mapsto(u+1/4,Dy,iz)$, and
$W=e^{-2\pi i u}z$ is invariant.  With core angle $a=4u\bmod1$
and normal angle $k=\arg W/(2\pi)$, the coarse angle is $a+4k$.
Let $g,h$ denote the two corresponding loops.  The integral basis
\[
 v=g^4h^{-1},\qquad t=g,
 \qquad [v,t]_{[g,h]}=\begin{pmatrix}4&1\\-1&0\end{pmatrix}
\]
has determinant one.  Its coarse degrees are zero and one; its
vertical monodromies are $T^4=I$ and $T$ respectively.
It specifies the generic six-torus fibre and the primitive
base loop in the boundary product.

\item A suspension cochain representing $[w]$ has values $(w,0)$
on $(g,h)$.  Following the four consecutive core prisms transports
their values by $I,T,T^2,T^3$, giving
\[
 (w,0)_{(g,h)}\longmapsto(N_Tw,w)_{(v,t)},
 \qquad N_T=I+T+T^2+T^3.
 \tag{FT.1}\hypertarget{tag-02J0}{}\label{maps:eq:FT.1}
\]
For a torsion coinvariant there are an integer $n\ne0$ and an
integral cochain $b$ with $nw=(T-I)b$.  Therefore
$nN_Tw=N_T(T-I)b=0$, and $N_Tw=0$ in the free cochain lattice.
The fibre restriction vanishes; its base-circle component is $[w]$.
Replacing $w$ by $w+(T-I)b$ changes that component by its cellular
coboundary.  The formula is consequently a well-defined map of the
stated integral coinvariant groups.

\item The base cover by a disc at the finite point and its
complement gives the relative \v{C}ech connecting map.  On their
annular intersection the degree-zero fibre restriction is zero and
the next component is the one-cocycle just computed.
The connecting map sends it to $\pm[w]$ in the global coinvariants;
the two ordinary peripheral loops impose the relations
$\bigwedge^{q-1}M_1^t-I$ and $\bigwedge^{q-1}M_2^t-I$.
These are the two successive boundaries in the same base--fibre
complex, so their composition is the Leray $d_2$ of
\eqref{maps:eq:F.6}.  The one boundary-orientation choice accounts
for the sign.  Its coefficient is one, as witnessed by the
degree-one loop $t$ above.  Point-supported kernels or cokernels in
the comparison of restriction sheaves have no positive cohomology,
so they preserve this degree-two base quotient.

\item In signed-permutation coordinates the negative orbits in
vertical degrees $1,2,3$ are represented respectively by
$e_4$, $e_0\wedge e_2$ and $e_0\wedge e_2\wedge e_4$.
The orbit relations give exactly one local $\mathbb Z/2$ in each
case.  Transport by $\bigwedge^kC_2^{-t}$ and then include vertically
in the six-lattice.  This is the same map as the vertical inclusion
followed by $\bigwedge^kQ_2$ used in
\cref{pub:V:integral-restrictions}, because $P_2=(-1)\oplus C_2$.
The explicit $w_k,\rho_k,E_k,z_k$ there satisfy
\[
 \rho_k\Delta_k=0\pmod2,\qquad
 \rho_kw_k=1\pmod2,\qquad E_kz_k=2w_k,
 \qquad |\det E_k|=2.
\]
Thus the unique nonzero local class maps to the unique nonzero
global class in each of the three degrees.  Formula
\eqref{maps:eq:FT.1} sends the zero representative to zero as well.
\end{enumerate}
\end{proof}




\begin{theorem}[Integral homology and the middle form]
\hypertarget{tag-02NW}{}\label{maps:src:F.2.5}
\hypertarget{tag-02NX}{}\label{maps:src:F.2.6}
\hypertarget{tag-02QL}{}\label{night:P1.12}
\StatementRecord{maps:src:F.2.6}{proved}{02NX}{}
The marked smooth gluing $M$ of \cref{pub:V:smooth-gluing} is
three-connected and has torsion-free integral homology with Betti numbers
\begin{equation}
 (1,0,0,0,4,0,0,0,1).
 \hypertarget{tag-02F9}{}\label{maps:eq:F.5}
\end{equation}
It admits an almost complex structure, has signature two, and its
middle integral intersection form is
$3\langle1\rangle\oplus\langle-1\rangle$.
The compact holomorphic realization of \cref{maps:src:F.3.7}
has this same marked smooth type.
\end{theorem}
\begin{proof}
The marked identification in \cref{maps:src:F.3.7} compares its
finite affine actions, cusp fan and overlap maps with the gluing
of \cref{pub:V:smooth-gluing}. The bounded-correction comparison
\cref{pub:V:cusp-smooth-dependence} preserves its collars.
We compute the integral groups of that gluing.
By \cref{pub:V:integral-restrictions} and
\cref{mapsaudit:order-four-torsion-transgression}, the point-supported
torsion classes have differential
\begin{equation}
 d_2[w]=\pm[w]\in L_{q-1}/\operatorname{im}\Delta_{q-1}.
 \hypertarget{tag-02FA}{}\label{maps:eq:F.6}
\end{equation}
The three images are the nonzero global classes, as shown by the
integer witnesses in \cref{pub:V:integral-restrictions}.
Since $H^0(\mathcal S_q)=0$ in these three degrees, the
point-supported kernels give
\[
 H^0(\mathcal R_2)=\mathbb Z/2,\quad
 H^0(\mathcal R_3)=\mathbb Z/2,\quad
 H^0(\mathcal R_4)=\mathbb Z\oplus\mathbb Z/2.
\]
The Leray sequence
\cite[Tag 01F2, Lemma 20.13.4; chapter PDF p.~22]{StacksProject}
has only base columns zero, one and two.
The maps from $(0,2)$ to $(2,1)$ and from $(0,3)$ to $(2,2)$ are
the isomorphisms just proved.  The restriction of the map from
$(0,4)$ to $(2,3)$ to its torsion summand is also an isomorphism.
Writing the image of a free cusp generator as $\epsilon\in\mathbb Z/2$,
its kernel is the free group generated by $(1,\epsilon)$.
All total-degree-three terms consequently vanish.  In total degree
four the remaining terms are $\mathbb Z$ and $\mathbb Z^3$.
Their extension is free, so
\[
                   H^3(M;\mathbb Z)=0,\qquad
                   H^4(M;\mathbb Z)=\mathbb Z^4.
\]
The degree-one row of the same Leray sequence is
\[
 H^0(\mathcal R_1)=\mathbb Z
 \xrightarrow{d_2^{0,1}}H^2(\mathcal R_0)=\mathbb Z.
\]
Its kernel is $H^1(M;\mathbb Z)=0$ by
\cref{pub:V:smooth-gluing}. Its degree is nonzero, so it becomes
an isomorphism over $\mathbb Q$. The image-sheaf groups in
\cref{pub:V:integral-restrictions} give
$H^2(M;\mathbb Q)=H^3(M;\mathbb Q)=0$.
Universal coefficients
\cite[Theorem~3.2 and Corollary~3.3, pp.~195--196]{Hatcher2002}
identify the torsion of $H_2$ and $H_3$ with the corresponding
$\operatorname{Ext}$ subgroups of $H^3$ and $H^4$; both are zero.
Their rational ranks also vanish, so $H_2=H_3=0$.
Integral Poincar\'e duality
\cite[Theorem~3.30, p.~241]{Hatcher2002} gives the remaining groups.
The successive Hurewicz homomorphisms
\cite[Theorem~4.32, p.~366]{Hatcher2002} give $\pi_2=\pi_3=0$.
Finally the almost complex structure and Euler number in
\cref{pub:V:smooth-gluing} satisfy
\cref{sec:V-known-inputs}, which gives the stated signature and
middle form.
\end{proof}






















\section{The oriented diffeomorphism type}

\begin{theorem}[Characteristic numbers and the middle lattice]\hypertarget{tag-02QM}{}\label{night:P1.13}
\StatementRecord{night:P1.13}{proved}{02QM}{\SeriesRef{night:P1.12}, \SeriesRef{sec:V-known-inputs}; \cite[equations (1.2)--(1.4), pp.~3--4]{ffEsoleKang}; \cite[Theorem 2.2.1, p.~61]{ffPetersSterk}}
With the complex orientation the fourfold has
\[
\sigma(M_c)=2,\qquad H^4(M_c;\mathbb Z)\cong I_{3,1},\qquad
w_2(M_c)=0,
\]
and there is a diagonal integral basis in which
$p_1/2=(1,1,1,1)$.
Its Chern numbers are
$c_4=6$, $c_2^2=2$, and
$c_1^4=c_1^2c_2=c_1c_3=0$; its Pontryagin numbers are
$p_1^2=8$, $p_2=14$.
\end{theorem}
\begin{proof}
\begin{enumerate}[label={$\langle1\rangle$\arabic*.}]
\item \Cref{night:P1.12} implies $c_1=c_3=0$ as integral cohomology classes;
$c_4[M_c]=\chi(M_c)=6$.  The Chern-root identities give
$p_1=-2c_2$, $p_2=c_2^2+2c_4$.
The symbol-index interpretation in \S\ref{sec:V-known-inputs},
with the polynomial \cite[equation (1.2), p.~3]{ffEsoleKang},
gives integrality of the Todd genus. Expanding
$\prod x_j/(1-e^{-x_j})$ to degree four gives
\[
\frac{3c_2^2-6}{720}\in\mathbb Z,
\quad\text{hence}\quad c_2^2=2+240k\quad(k\in\mathbb Z).
\]
The signature polynomial \cite[equation (1.4), p.~4]{ffEsoleKang}, applied to
this closed oriented smooth eight-manifold gives
\[
\sigma=\frac{7p_2-p_1^2}{45}
      =\frac{3c_2^2+84}{45}=2+16k.
\]
The rank-four middle form has $|\sigma|\leq4$, so $k=0$.
This proves the displayed characteristic numbers.

\item Duality makes the middle form unimodular.  Its signature is two and
its positive and negative indices are three and one.  The rank-at-most-four integral classification
\cite[Theorem 2.2.1, p.~61]{ffPetersSterk} applied to these
positive and negative indices gives the odd form \(I_{3,1}\).
Since $H^2(M_c;\mathbb Z/2)=0$, its $w_2$ is zero.
For the complex tangent bundle,
$w_4=c_2\bmod2=(p_1/2)\bmod2$.
The Wu identity in degree four, with $w_1=w_2=0$, says that $w_4$ is
characteristic for the middle pairing.  Thus the vector $p_1/2$ is
characteristic and has square $2$.

\item We give the required integral isometry of this pair explicitly as
a terminating procedure.  In a diagonal basis a characteristic vector
has four odd coordinates.  Change their signs to write it as
$(a,b,c,d)$ with all positive and
$a^2+b^2+c^2=d^2+2$.
For $d>1$,
\[
d<a+b+c<3d.
\]
The first inequality follows by squaring and using the positive cross
terms; the second follows from Cauchy--Schwarz and $d>1$.
Reflection in the integral norm-two vector $(1,1,1,1)$ sends this
vector to itself minus $(a+b+c-d)(1,1,1,1)$.
Its fourth coordinate is $2d-a-b-c$, whose absolute value is strictly
smaller than $d$.  The reflection preserves the integral lattice and
characteristic parity.  Repeating, with sign changes at each step,
terminates at $d=1$, when the norm equation forces $a=b=c=1$.
This supplies an integral isometry taking $p_1/2$ to $(1,1,1,1)$,
and proves the claimed pair, not just its square.

\end{enumerate}
\end{proof}




\begin{theorem}[a complex structure on the stated connected sum]\hypertarget{tag-02OC}{}\label{maps:src:F.4.1}
\StatementRecord{maps:src:F.4.1}{proved}{02OC}{\SeriesRef{maps:src:F.3.7}, \SeriesRef{maps:src:F.2.6}, \SeriesRef{sec:V-known-inputs}, \SeriesRef{night:P1.13}; \cite[Corollary 1.4, p.~4]{ffMandal}; \cite[Definition 3.6 and Theorem 3.7(1), pp.~12--14]{IVVWallForm}; \cite[Theorem A, p.~178]{IVVKramerStolz}}


The compact complex fourfold of Theorem \ref{maps:src:F.3.7} is three-connected,
has Euler number six, and is orientation-preservingly diffeomorphic to
\[
 3\mathbb H\mathbb P^2\ \#\ (-\mathbb H\mathbb P^2).
 \tag{F.25}\hypertarget{tag-02EW}{}\label{maps:eq:F.25}
\]
Its integral middle lattice is $I_{3,1}$ and, in a diagonal integral
basis, $p_1/2=(1,1,1,1)$.  

\end{theorem}


\begin{proof}
\begin{enumerate}[label={$\langle1\rangle$\arabic*.}]
\item The integral calculation in \cref{maps:src:F.2.6} gives
$\pi_1(M)=0$, $H_2(M;\mathbb Z)=H_3(M;\mathbb Z)=0$, and
torsion-free homology of ranks $1,4,1$ in degrees $0,4,8$.
Hurewicz applied successively gives $\pi_2(M)=\pi_3(M)=0$, and
$e(M)=6$. The complex orientation is the one in
\cref{maps:src:F.3.7}. The characteristic calculations of \cref{night:P1.13} give
\[
Q=I_{3,1},\qquad c_2^2=2,\qquad
p:=p_1/2=-c_2,\qquad p^2=2.
\]
Here $c_1=c_3=0$ because the corresponding integral groups vanish.
The spin Wu identity makes $p$ characteristic for $Q$. Evaluating
its characteristic congruence on each diagonal basis vector shows
that all four coordinates of $p$ are odd.

\item Put $v=(1,1,1,1)^t$.  Since $v^tQv=2$, the reflection $R=I-v(v^tQ)$ is an integral isometry of $Q$ with $R^2=I$.  By the descent in the proof of \cref{night:P1.13}, a product $U$ of such reflections and of coordinate sign changes satisfies $Up=v$.

\item We compute the complete normal clutching datum in fixed
coordinates. Orient $\mathbb H$ by $(1,i,j,k)$ and use
\[
\operatorname{Spin}(4)=S^3_L\times S^3_R,
\qquad (u,w)h=uhw^{-1}.
\]
The covering induces an isomorphism on $\pi_3$. Let $(a,b)$ be
the degrees of the two clutching maps. Evaluation on $1$ gives
$uw^{-1}$, of degree $a-b$, so $e=a-b$, with the sphere orientation
fixed by the left degree-one Euler generator. The left action is
complex linear for right multiplication by $i$, whose real complex
basis $(1,i,j,-k)$ has the opposite quaternion orientation. Its
degree-one generator therefore has $c_2=-1$, and hence $p_1=2$.
The right inverse action is complex linear for left multiplication
by $i$, with real basis $(1,i,j,k)$. Its Euler number and $c_2$ are
both $-1$, so its $p_1$ is again $2$. The determinant-one complex
actions have $c_1=0$. Consequently the exact homomorphism and its
inverse on its image are
\[
(a,b)\longmapsto(e,p_1/2)=(a-b,a+b),
\qquad
(a,b)=\left(\frac{e+p_1/2}{2},\frac{p_1/2-e}{2}\right).
\]
The image consists of pairs of the same parity, and the map is
injective. The tangent clutching class of $S^4$ has $(e,p_1)=(2,0)$,
hence coordinates $(1,-1)$. The homotopy sequence of
$SO(4)\to SO(5)\to S^4$ identifies its span with the kernel of
stabilization and gives surjectivity on $\pi_3$. Subsequent
stabilizations are isomorphisms there, since
$\pi_3(S^n)=\pi_4(S^n)=0$ for $n\geq5$.
The stable coordinate is therefore $a+b=p_1/2$, with the left
generator sent to $1$.

\item For $x\in H_4(M;\mathbb Z)$, use the embedded sphere
representative in the setup of
\cite[\S3.3, pp.~11--12]{IVVWallForm}. Orient its normal bundle
$\nu_x$ by the ambient orientation. Self-intersection and the
tangent splitting give
\[
e(\nu_x)=x\cdot x,\qquad
\langle p_1(\nu_x),[S^4]\rangle=\langle p_1(TM),x\rangle,
\]
because $TS^4$ is stably trivial. Thus its two clutching coordinates
are exactly
\[
\left(\frac{x\cdot x+\langle p,x\rangle}{2},
      \frac{\langle p,x\rangle-x\cdot x}{2}\right),
\qquad S\alpha_M(x)=\langle p,x\rangle.
\]
Their integrality is the characteristic congruence in step~1.
This is the stable normal homomorphism of
\cite[Definition~3.6, p.~12]{IVVWallForm} with its generator.
For an oriented middle-cohomology isometry $\psi$ from $M$ to a
comparison manifold $N$, the corresponding homology map is
\[
\Phi=\operatorname{PD}_N\circ\psi\circ\operatorname{PD}_M^{-1}.
\]
When $\psi(p_M)=p_N$, the cup-product evaluation identity gives
$\langle p_N,\Phi(x)\rangle=\langle p_M,x\rangle$.
The map $\Phi$ preserves the intersection pairing and both displayed
normal coordinates, so it intertwines the complete extended forms.

\item Take $N=3\mathbb H\mathbb P^2\#(-\mathbb H\mathbb P^2)$.
On each summand choose the generator $u$ with $p_1=2u$, as in
\cite[Lemma~2.1, p.~5]{ffMandal}. Reversing orientation changes
the evaluation of $u^2$ and leaves the tangent Pontryagin class
unchanged. The collapse maps of the oriented connected sum give
\[
H^4(N;\mathbb Z)=\bigoplus_{i=1}^4\mathbb Z u_i,
\qquad Q_N=I_{3,1},\qquad p_N=u_1+u_2+u_3+u_4;
\]
the tangent formula is
\cite[Theorem~2.4, p.~7]{ffMandal}.
Van Kampen, the connected-sum homology sequence and Hurewicz show
that $N$ is three-connected. The isometry $U$ of step~2 is an
isometry $\psi:H^4(M;\mathbb Z)\to H^4(N;\mathbb Z)$ carrying
$p_M$ to $p_N$. Step~4 supplies the isomorphism of their smooth
extended forms. Wall's classification in
\cite[Theorem~3.7(1), pp.~13--14]{IVVWallForm} now gives an
orientation-preserving diffeomorphism
$M\cong N\#\Sigma$ for an oriented homotopy eight-sphere $\Sigma$.

\item Kramer--Stolz's
\cite[Theorem~A, p.~178]{IVVKramerStolz} applies to
$\mathbb H\mathbb P^2$: it is smooth, simply connected, has dimension
eight, and has total integral homology $\mathbb Z\oplus\mathbb Z
\oplus\mathbb Z$. It gives
$\mathbb H\mathbb P^2\#\Sigma\cong\mathbb H\mathbb P^2$.
This diffeomorphism preserves orientation: its map on $H^4$ sends
a generator to that generator or its negative, whose square still
evaluates to $1$, forcing its degree to be $+1$.
The oriented associativity and permutation maps for connected sum
place $\Sigma$ beside one positive summand. Choose the punctured
balls compatibly with that diffeomorphism and use the identity on
the remaining punctured summands. These maps give
\[
\begin{aligned}
N\#\Sigma
&\cong(\mathbb H\mathbb P^2\#\Sigma)\#\mathbb H\mathbb P^2
       \#\mathbb H\mathbb P^2\#(-\mathbb H\mathbb P^2)\\
&\cong N.
\end{aligned}
\]
Composing with step~5 supplies the required oriented smooth
diffeomorphism. Transporting the integrable structure of
\cref{maps:src:F.3.7} along it gives the asserted complex structure.
\end{enumerate}
\end{proof}


\section{The normalized connection and the exact parameter domain}










\begin{theorem}[Certified values of $h_b$]\hypertarget{tag-02QP}{}\label{night:P1.16}
\StatementRecord{night:P1.16}{proved}{02QP}{\SeriesRef{night:P1.15}, \SeriesRef{night:P1.18}, \SeriesRef{night:P1.8}}
In the canonical splitting $\sigma_b$ of \cref{night:P1.15},
\[
\begin{aligned}
-1.0342766224163217994942488453783461&<h_b(1),\\
h_b(1)&<-1.0342766224163217994942488453783460.
\end{aligned}
\]
In particular $h_b(1)<-1$.  The parameter $c=5/(3\pi)$, the exceptional value of \cref{night:P1A.2}, lies outside the global transversality domain.  
\end{theorem}
\begin{proof}
\begin{enumerate}[label={$\langle1\rangle$\arabic*.}]
\item \textbf{The rational connection and the marked paths.}
Put $a=256/3125=-\lambda_*$ and
\[
 D(\lambda)=\operatorname{diag}(1,\lambda^2,\lambda,1,
                              \lambda^{-1},\lambda^{-2}).
\]
For the row-frame equation $Y'=Y\Gamma_6$ of \cref{night:P1.15}, the exact
change $\widetilde Y=YD$ gives
\[
 \widetilde\Gamma
 =D^{-1}\Gamma_6D+D^{-1}D'
 =\frac{R_0}{\lambda}+\frac{R_*}{\lambda+a},\qquad
 R_0=\lim_{\lambda\to0}\lambda\widetilde\Gamma,
 \quad R_*=\lim_{\lambda\to-a}(\lambda+a)\widetilde\Gamma.
\]
These formulas specify both rational residue matrices entry by entry
from (1.15.1)--(1.15.2).  Direct rational multiplication verifies the
identity.  Their maximum absolute row-sum norms are
\[
 \|R_0\|_\infty=\frac{3125}{256},\qquad
 \|R_*\|_\infty=\frac{14575}{768}.
\]
Let $r=1/32$.  The zero meridian is the polygon with successive
vertices
\[
 1,\ r,\ r+ir,\ -r+ir,\ -r-ir,\ r-ir,\ r,\ 1.
\]
The infinity meridian has successive vertices
\[
 1,\ 1-i,\ -1-i,\ -1+i,\ 1+i,\ 1.
\]
The first runs counterclockwise about $0$ and excludes $-a$; the
second runs clockwise about both finite poles, hence positively
about infinity.  The zero square has radial distance between $r$
and $\sqrt2r<a$, and the infinity square between $1$ and $\sqrt2$,
with $1>a$.  Linear interpolation of these positive radial distances
to $r$ and $1$, respectively, gives based homotopies through the
corresponding punctured regions to their circular meridians; the
zero stem remains the positive segment.  Thus these are the based meridians of \cref{night:P1.15}, with matrices $B$ and $A$, respectively.
Since $D(1)=I$, the gauge
keeps their based matrices unchanged.  The open comparison path
from $1$ to $2$ is the positive real segment.

\item \textbf{The Taylor recurrence and its complete error bound.}
Write $M_1=R_0+R_*$ and $M_0=aR_0$.  At a rational complex center
$z$, set
\[
 Q_0=z(z+a),\quad Q_1=2z+a,\quad P_0=zM_1+M_0.
\]
For the fundamental solution $T(t)=\sum_{n\ge0}T_nt^n$ normalized
by $T_0=I$, multiplication of
$T'=T(M_1(z+t)+M_0)/((z+t)(z+a+t))$ gives, with $T_{-1}=0$,
\[
 (n+1)Q_0T_{n+1}
 =T_nP_0+T_{n-1}M_1-nQ_1T_n-(n-1)T_{n-1}.
\]
This recurrence uses only rational connection coefficients and
outward-rounded complex balls.

For each step put
\[
 d_0=\max(|\operatorname{Re}z|,|\operatorname{Im}z|),\quad
 d_*=\max(|\operatorname{Re}(z+a)|,|\operatorname{Im}z|),
 \quad R=\tfrac12\min(d_0,d_*).
\]
The segment is subdivided with rational increments $h$ satisfying
$|h|_1=|\operatorname{Re}h|+|\operatorname{Im}h|\le R/4$.
The whole complex disc $|t|\le R$ avoids both poles, and on it
\[
 \|\widetilde\Gamma(z+t)\|_\infty\le
 M=\frac{\|R_0\|_\infty}{d_0-R}
   +\frac{\|R_*\|_\infty}{d_*-R}.
\]
The integral equation and Gronwall's estimate
\cite[Lemma~2.7, pp.~42--43]{TeschlODE2012}, followed by Cauchy's
coefficient formula
\cite[\S1.9(iii), (1.9.30)--(1.9.31)]{DLMFCauchy2026}, give the
entrywise degree-$N$ tail bound
\[
 E_N=\exp(MR)\frac{(|h|_1/R)^{N+1}}{1-|h|_1/R}.
 \tag{1.16.1}
\]
Adding the real and imaginary error balls of radius $E_N$ to each
entry therefore encloses the exact transition.  All subsequent
matrix products and inverses retain those radii.  The order starts
at $32$ and is increased by $16$ until $E_N<10^{-100}$.
There are $120$ steps for the zero meridian, $66$ for the infinity
meridian and $6$ for the open path, hence $192$ in total; the
largest order is $192$.  

\item \textbf{The normalized intertwiner and the height.}
Let $H_0,H_\infty$ denote the two computed based transition matrices.
The exact equations are
\[
 XH_0=BX,\qquad XH_\infty=AX,\qquad Xe_0=e_0.
\]
\Cref{night:P1.15} proves their simultaneous existence and uniqueness.  The
fixed first column leaves thirty unknowns.  We select thirty of
the sixty scalar equations in the other five columns and certify
their coefficient matrix inverse in complex-ball arithmetic.
The exact common solution therefore belongs to the resulting
matrix ball; the remaining equations are also checked to contain
zero.  Their consistency is given by \cref{night:P1.15}.

Likewise, the real intersection in (1.15.3) is computed from an
invertible four-equation subsystem in
$(\operatorname{Re}v_2,\operatorname{Im}v_2,
  \operatorname{Re}v_1,\operatorname{Im}v_1)$.
\Cref{night:P1.8} proves the existence and uniqueness of the full intersection;
the unused equation is checked with its interval radius retained.
At $\lambda=1$ the resulting height ball has radius less than
$4\cdot10^{-91}$.  It is contained in the following wider displayed
interval:
\[
\begin{aligned}
 h_b(1)\in\bigl[&
 -1.03427662241632179949424884537834605259496477790218821132\\
 &\qquad\mathbin{\pm}8.11\cdot10^{-57}\bigr].
\end{aligned}
\]
It lies strictly between the two decimal endpoints in the statement.
Since $5/(3\pi)<1$, we have $h_b(1)+5/(3\pi)<0$.
The cusp asymptotic in \cref{night:P1.8} gives $h_b+c\to+\infty$ along a path
to the cusp.  Continuity on a path from $1$ to that region therefore
supplies a zero for $c=5/(3\pi)$, proving its inadmissibility.

\item \textbf{The regular-root evaluation and the exact comparison.}
The second method uses the regular frames with powers
$(0,6,3,0,-3,-6)$ for $\lambda=u^4$ at zero and
$(0,-8,-4,0,4,8)$ for $\lambda=u^{-5}$ at infinity.
Their transformed connections have denominators
$3125u^4+256$ and $3125+256u^5$.  The regular-root recurrences,
finite deck characters and normalized intertwiner construction
are displayed in the proof of \cref{night:P1.18}.  They specify $X_{\rm root}$
and the constant infinity frame $F_\infty$ without integrating a
closed loop.

Let $T_{12}$ be the Fuchsian transition on the positive segment.
The first method gives
\[
 Y_2=XT_{12}D(2)^{-1}.
\]
For the second method put $u=2^{-1/5}$ and
$D_\infty(u)=\operatorname{diag}(1,u^{-8},u^{-4},1,u^4,u^8)$.
Let $U_\infty$ solve $U_\infty'=U_\infty\Gamma_\infty$ with $U_\infty(0)=I$.  The second evaluation is
\[
 Y_{2,\rm root}=F_\infty U_\infty(2^{-1/5})
                         D_\infty(2^{-1/5})^{-1}.
\]
The explicit polynomial numerator of $\Gamma_\infty$ is obtained
by the stated gauge and chain rule from $\Gamma_6$; its coefficient
matrices are saved.  On $|u|\le5/4$, its denominator has modulus
at least $2343.75$, and the row-sum connection bound is $65.6375$.
The degree-$880$ Taylor polynomial and (1.16.1) consequently give
an error less than $5.39\cdot10^{-103}$ at $2^{-1/5}$.
The complete ball comparisons give
\[
 \max_{i,j}|(X-X_{\rm root})_{ij}|<1.37\cdot10^{-73},\qquad
 \max_{i,j}|(Y_2-Y_{2,\rm root})_{ij}|<4.71\cdot10^{-74}.
\]
Thus both period matrices agree to more than thirty decimal digits
under the explicit frame maps above.

\item \textbf{The period planes.}
For the real parameter $c$ the period matrix is
\[
 W_c=(Y_{*,1},\,Y_{*,2},\,Y_{*,3}-(625\pi/96)c\,e_0).
\]
Conjugation of $\det(W_c,\overline{W_c})$ interchanges two blocks
of three columns, so it multiplies the determinant by $(-1)^9=-1$.
Thus the determinant is purely imaginary.  At the two points above
and $c=10$, its positive imaginary parts are
\[
\begin{gathered}
 88279007523.2167451846066607088299709479845\ldots,\\
 3178499608.3637608517583236010752861132441\ldots .
\end{gathered}
\]
The complete stored balls have radii less than $4\cdot10^{-79}$
and $2\cdot10^{-80}$, respectively.  These are evaluations at the
two specified points.  The global admissibility of $c=10$ follows
from the complete minimum argument in \cref{night:P1.18}.  The regular-root
calculation there also gives
\[
\begin{split}
 h_b(\infty)&=-1.094540250311402780116268263054945703725533\ldots,\\
 h_b(0)&=-0.931793376356523505430795105796011114990828\ldots .
\end{split}
\]

The computation is done by the program \texttt{period\_loop\_comparison.py} in interval arithmetic (python-flint, Arb) with 140 decimal digits; a second program, \texttt{validate\_saved\_period\_comparison.py}, rechecks the saved enclosures and every displayed bound.
\end{enumerate}
\end{proof}




\begin{theorem}[The minimum of the height function]
\hypertarget{tag-02QR}{}\label{night:P1.18}
\StatementRecord{night:P1.18}{proved}{02QR}{\SeriesRef{night:P1.15}, \SeriesRef{night:P1.8}}
The function $h_b$ has its unique minimum at $\lambda=\infty$.
For the normalized fundamental solution $H(1)=I$ of the algebraic
connection and its regular infinity limit $H_\infty$, this minimum is
\[
 h_{b,\min}=-\frac{96}{625\pi}\operatorname{Re}(XH_\infty)_{03}.
\]
Its negative $c_*=-h_{b,\min}$ satisfies the exact bounds
\[
 1.094540250311402780116268263054945703725<c_*
 <1.094540250311402780116268263054945703726.
 \hypertarget{tag-03KZ}{}\label{layout:V:fresh-bounds}
\]
The full parameter domain is $\mathcal U=\{\xi:\operatorname{Im}\xi>c_*\}$.
\end{theorem}
\begin{proof}
Put $a=256/3125$ and $c_0=96/(625\pi)$. We use the rational
connection of \cref{night:P1.15}, its integral matrices $A,B$,
and the normalization $Y=XH$, $H(1)=I$, $Xe_0=e_0$.
The finite-root gauges are
\[
 \lambda=u^4,\quad D_0=\operatorname{diag}(1,u^6,u^3,1,u^{-3},u^{-6}),
 \qquad
 \lambda=u^{-5},\quad D_\infty=
 \operatorname{diag}(1,u^{-8},u^{-4},1,u^4,u^8).
 \hypertarget{tag-03KX}{}\label{layout:V:fresh-gauges}
\]
For either gauge the regular connection is
$D^{-1}\lambda'\Gamma_6D+D^{-1}D'$.
Let $U_0,U_\infty$ be its solutions with value $I$ at zero.
The deck characters $T_0,T_\infty$ give the exact based
monodromies $U_0(1)^{-1}T_0^{-1}U_0(1)$ and
$U_\infty(1)^{-1}T_\infty^{-1}U_\infty(1)$.
The unique common intertwiner fixing $e_0$ is $X$, by
\cref{night:P1.15}. Hence the regular limiting frames are
$F_0=XU_0(1)^{-1}$ and $F_\infty=XU_\infty(1)^{-1}$.

The height is obtained from a real linear system. Write
$q=(0,v_2,v_1,-i,0,0)^t$ and impose that the lower five
coordinates of $Yq$ be real. The four real unknowns are the
real and imaginary parts of $v_2,v_1$; existence and uniqueness
are those of \cref{def:hb}. In these coordinates
$h_b=c_0\operatorname{Im}(Yq)_0$.
Every monodromy is real and fixes $e_0$, so the imaginary top
coordinate is unchanged when the lower coordinates are real.
This proves that the calculation gives one function on the base.
The real normalized conjugation matrix $K=X\bar X^{-1}$ satisfies
$Y(\bar\lambda)=K\overline{Y(\lambda)}$ and $Ke_0=e_0$;
using $-\bar q$ gives $h_b(\bar\lambda)=h_b(\lambda)$.

We give the enclosure method and the resulting whole-domain
inequalities. For an analytic right connection $G$ on a disc
$|z-z_0|\le R$, with a submultiplicative norm bound $\|G\|\le M$,
its normalized solution satisfies $\|U\|\le e^{MR}$.
Gronwall's estimate and Cauchy's coefficient formula therefore
bound the tail after degree $N$, on $|z-z_0|\le r<R$, by
\[
 e^{MR}\frac{(r/R)^{N+1}}{1-r/R}.
\]
For rational $G=P/Q$, the positive lower bound
$|Q_0|-\sum_{j\ge1}|Q_j|R^j$ gives $M$ directly.
These estimates, with outward interval operations, enclose the
root transports, the intertwiner and the four-equation real
system. Existence of the exact intertwiner is supplied by
\cref{night:P1.15}; an invertible square subsystem encloses it.
A matrix inverse on a rectangle and differentiation of
$Mv=b$ give
\[
 v_s=M^{-1}(b_s-M_sv),\qquad
 v_{st}=M^{-1}(b_{st}-M_{st}v-M_sv_t-M_tv_s).
\]
Together with $Y_x=YG$, $Y_y=iYG$ and
$Y_{xx}=Y(G^2+G')$, these enclose the real Hessian of $h_b$
on whole rectangles. For the value on a rectangle of
half-diagonal $\rho$, an entrywise bound $|G|\le\mathcal A$
gives
$|Y(z)-Y(z_0)|\le|Y(z_0)|_+(e^{\rho\mathcal A}-I)$.
All inputs to these enclosures are the displayed rational
connection, its gauges, the integral monodromies and the specified
normalization.

At infinity the deck action sends the normalized coefficient
pair to $(\zeta^8v_2,\zeta^4v_1)$, with
$\zeta=e^{2\pi i/5}$. Its value at the root point is therefore
zero. The height there is
$h_b(\infty)=-c_0\operatorname{Re}(F_\infty)_{03}$,
and the same deck action makes its gradient zero.
In the scaled coordinate $w=a^{1/5}u$, so $\lambda=aw^{-5}$,
the enclosures on the closed square $[-1/10,1/10]^2$ give
\[
 \min(h_{xx},h_{yy})-|h_{xy}|>\frac14.
\]
The real Hessian is consequently at least $I_2/4$ there.
Integrating along the segment from zero gives
\[
 h_b(aw^{-5})-h_b(\infty)\ge\frac18|w|^2.
\]
The normalized endpoint enclosure has radius below $4\cdot10^{-76}$
and lies between the negatives of the two endpoints in the statement.
\hypertarget{tag-03KY}{}\label{layout:V:fresh-minimum}

For the rest of the base use $v=a^{-1/4}u_0$, so
$\lambda=av^4$, and put $R=49/50$, $S=103/100$.
The first-quadrant $w$-disc $|w|\le R$ covers
$|\lambda|\ge aR^{-5}$ under its finite deck maps; the
first-quadrant $v$-disc covers $|\lambda|\le aR^4$.
The collar $R\le|w|\le S$ covers the intervening annulus,
since $R^4\ge S^{-5}$.
The cusp coordinate is $q=\lambda/\lambda_*-1=-w^{-5}-1$.
The same enclosure formulas on rational rectangular partitions
give the following strict lower bounds away from the central
square and the inner cusp disc:
\[
\begin{array}{c|c}
\text{region}&\text{lower bound for }h_b\\\hline
\text{infinity-root disc}&-1.094216\\
\text{zero-root disc}&-1.083056\\
\text{cusp annulus }1/100\le|q|\le1/4&-1.083777\\
\text{collar}&-1.030556
\end{array}
\]
Every endpoint in this table is an exact terminating decimal
strictly greater than $h_b(\infty)$. Rectangles meeting a
boundary are kept closed. Those lying in another of the stated
regions are handled there; those outside a specified radius are
excluded by their exact squared-distance inequality. The deck
maps and conjugation give the remaining sectors. Thus the
bounds cover all points outside $|q|<1/100$ and the central
square.

For the inner cusp disc, the exact gauge is
$D_\lambda=\operatorname{diag}(1,\lambda^2,\lambda,1,\lambda^{-1},\lambda^{-2})$.
Its connection is $R_c/q+A_0/(1+q)$ with $R_c^2=0$.
Writing $YD_\lambda=H_c(I+R_c\log q)P(q)$, $P(0)=I$,
the coefficient recurrence is
\[
 (nI-\operatorname{ad})P_n
 =\sum_{j<n}(-1)^{n-1-j}P_jA_0,
 \qquad\operatorname{ad}(Q)=QR_c-R_cQ.
\]
Here $\operatorname{ad}^3=0$, so the inverse is
$n^{-1}(I+\operatorname{ad}/n+\operatorname{ad}^2/n^2)$.
A scalar majorant $(1-|q|)^{-\gamma}$ gives the same
Cauchy tail control. Continuation to the positive cusp anchor
fixes $H_c$ and the logarithm branch. In the unimodular marking
$J$ with $J^{-1}NJ=\left(\begin{smallmatrix}0&I\\0&0\end{smallmatrix}\right)$,
the lower graph block is invertible throughout the closed
inner disc. The marking used for this anchor and the initial
cusp marking are related by $N=AN'A^{-1}$; their change has
block form $\left(\begin{smallmatrix}D&E\\0&D\end{smallmatrix}\right)$
and sends $Z$ to $DZD^{-1}+ED^{-1}$. It preserves the
real-intersection criterion.
\hypertarget{tag-03L0}{}\label{layout:V:fresh-marking}

At the parameter $\xi=ic_*$ write the resulting graph as
$Z(q)=\log q\,I_3/(2\pi i)+C(q)$, and put
$y_0=\log100/(2\pi)$ and $t=-\log|q|/(2\pi)-y_0\ge0$.
The coefficients, in increasing degree, of
$\det(tI_3+y_0I_3+\operatorname{Im}C(q))$ are bounded below
on $|q|\le1/100$ by
$3/2,9/2,15/4,1$, respectively.
The determinant is therefore positive for every nonzero $q$
in that disc. A graph vector $(Zv,v)$ is real precisely when
$v$ is real and $\operatorname{Im}Z\,v=0$. Thus the
$c_*$ plane is transverse there. The height at $q=1/100$
lies between $-0.451$ and $-0.450$; connectedness and
\eqref{maps:eq:F.22} give
$h_b(q)>h_b(\infty)$ on the entire punctured inner disc.

The strict central Hessian inequality, the four outer bounds and
the inner cusp inequality prove that infinity is the unique
minimum. The exact domain now follows from
\cref{night:P1.8}.
\hypertarget{tag-03L1}{}\label{layout:V:fresh-splitting}
\end{proof}


\section{Cohomology, deformations and isomorphisms}

We keep the complex orientation and the integral marking of \cref{night:P1.1}.
Let $\widehat B$ be the root orbifold of orders $5,4$, let
$h:M\to\widehat B$ be the map given by the two multiple divisors,
and let $\pi:\widehat B\to\mathbb P^1$ be its coarse map.
Write $L=T^2$; in the pencil coordinate it is
\[
 L=\mathcal O_{\widehat B}\bigl(\tfrac45[\infty]
                  -\tfrac34[0]\bigr),\qquad \deg L=\tfrac1{20}.
\]
The cusp point has line bundle $L^{20}$, and
$T_{\widehat B}=L^9$. These identities include the stabilizer
characters: their quotients have integral divisors of degree zero
on $\mathbb P^1$. Denote the cusp by $W$, its normalization by $Y$,
and its reduced singular locus by $D$.

\begin{theorem}[Direct images and the canonical bundle]\hypertarget{tag-02R1}{}\label{night:P1A.1}
\StatementRecord{night:P1A.1}{proved}{02R1}{\SeriesRef{night:P1.10}, \SeriesRef{night:P1.8}, \SeriesRef{night:P1.9}}
For every admissible member, its vertical translation bundle on
$\widehat B$ and its first coherent direct image are
\[
 E_0=F^\vee=\mathcal O\oplus L^{-1}\oplus L^{-2},\qquad
 G=R^1h_*\mathcal O_M=\mathcal O\oplus L^{-1}\oplus L^{-2}.
\]
Moreover $R^qh_*\mathcal O_M=\bigwedge^qG$ for $0\le q\le3$,
$K_M=h^*L^{-6}$, and
\[
 (h^0(\mathcal O_M),\ldots,h^4(\mathcal O_M))=(1,1,0,0,0).
\]
\end{theorem}
\begin{proof}
\begin{enumerate}[label={$\langle1\rangle$\arabic*.}]
\item The evaluation of the marked lattice on the lifted plane $F$
identifies the Lie bundle of the smooth tori with $F^\vee$.
\Cref{night:P1.8} gives $F=\mathcal O\oplus L\oplus L^2$; its lift as a bundle
has the same splitting. The logarithmic toric vector fields extend
this Lie bundle over the cusp, so the identification holds on the
whole orbifold, including both finite root charts.

\item On the smooth locus the constant-cocycle description of torus
cohomology gives $G=V/F$. In the rank-five orthogonal quotient,
$K=\ker(\alpha:P\to\mathcal O)=L\oplus L^2$ is the radical of
the restricted pairing on $P$. Indeed every wedge containing the
fixed line pairs to zero with another such wedge except through
their contraction components; removing that scalar component
leaves exactly $K$. The nondegenerate pairing consequently
identifies $W/P$ with $K^\vee$. Thus
\[
 0\longrightarrow\mathcal O e\longrightarrow V/F
   \longrightarrow L^{-1}\oplus L^{-2}\longrightarrow0.
\]
For every integer $j$, descent by invariant local Laurent series gives
\begin{equation}\hypertarget{tag-01O3}{}\label{eq:P1Adegree}
 \pi_*L^j=\mathcal O_{\mathbb P^1}(d(j)),\qquad
 d(j)=\lfloor4j/5\rfloor+\lfloor-3j/4\rfloor.
\end{equation}
The extension group $H^1(L\oplus L^2)$ therefore vanishes, so this
sequence splits as a holomorphic vector-bundle sequence.

\item We verify that this quotient is the coherent extension at the cusp.
The normalization resolution of $\mathcal O_W$ has one toric
threefold, seven toric surfaces, twelve toric curves and six points.
All positive coherent cohomology of those smooth complete toric
strata vanishes \cite[Theorem 9.2.3, p.~410]{CoxLittleSchenck2011}. Their constants form the cellular cochain complex of the periodic
Kuhn triangulation modulo $\mathbb Z^3$: its simplices are the faces of
$(0,e_{\pi(1)},e_{\pi(1)}+e_{\pi(2)},(1,1,1))$ for
$\pi\in S_3$, identified by integral translations. The quotient map
$x\mapsto x\bmod\mathbb Z^3$ identifies its realization with the dual
three-torus and respects every face map. Its cohomology is
$\bigwedge^q\mathbb C^3$, with the exterior cup product.
The face incidences, unit Smith factors and the Alexander--Whitney
products of the three coordinate cocycles are reconstructed in
\href{https://github.com/KokunoYumeto/s6-complex-structure-workbench/blob/83d76022bcd442278b760db9ff438afceed6c6ea/pro_returns/S6_SERIES_GRIND_checkpoint_10_20261120T0211Z/code/current_audits/audit_corrections/code/independent_V_analytic_06.py}{\texttt{independent\_V\_analytic\_06.py}} and its JSON output.
The invariant flat cohomology $\ker N$ maps to these classes by
constant Cech cocycles on the toroidal covering. A component
translation in that cover evaluates the corresponding lattice
coordinate. In the nilpotent logarithmic frame, $\ker N$ maps
isomorphically to $V/F$, since $N:F\to\ker N$ is an isomorphism.
Thus these are the extending quotient classes.

\item Flatness over each smooth root disc follows because its base is a
DVR and the smooth total space is torsion free over it. Properness was proved in \cref{night:P1.9,night:P1.10}. The coherent fibre dimensions are the
constant binomial dimensions on smooth tori, finite root fibres,
and the cusp. \cite[\S7, Satz~5]{Grauert60}, with these hypotheses,
gives local freeness and base change. The cup-product maps are
isomorphisms on every fibre, including the preceding central
cellular model, so they identify the remaining direct images with
$\bigwedge^qG$.

\item There is an exact sequence of locally free logarithmic bundles
\begin{equation}\hypertarget{tag-01O4}{}\label{eq:P1Alogcotangent}
0\longrightarrow h^*L^{11}\longrightarrow\Omega_M^1(\log W)
 \longrightarrow h^*F\longrightarrow0.
\end{equation}
In a chart $q=x_0\cdots x_r$, its first arrow sends
$dq/q$ to $\sum dx_i/x_i$ and its quotient is the height-zero
logarithmic cotangent bundle of the torus. These formulas also
prove exactness on every stratum. Taking determinants and
cancelling $\mathcal O_M(W)=h^*L^{20}$ yields
$K_M=h^*L^{11+3-20}=h^*L^{-6}$.

\item For any bundle which is a sum of powers of $L$, projection and
the exact finite averaging projector compute its direct images
by tensoring with $\bigwedge^qG$ and applying
\eqref{eq:P1Adegree}. All exponents $-1,\ldots,-5$ have coarse
degree $-1$, while exponent zero has degree zero. The Leray
sequence on the curve then gives the asserted structure-sheaf
cohomology. This includes the nonreduced coarse fibres, since
the computation was made before exact finite descent.

\end{enumerate}
\end{proof}




\begin{theorem}[The Hodge numbers]\hypertarget{tag-02R2}{}\label{night:P1A.2}
\StatementRecord{night:P1A.2}{proved}{02R2}{\SeriesRef{night:P1.15}, \SeriesRef{night:P1A.1}}
The complete Hodge table of every admissible fourfold is the
following, with rows indexed by $p$ and columns by $q$:
\[
\begin{array}{c|ccccc}
h^{p,q}&0&1&2&3&4\\\hline
0&1&1&0&0&0\\
1&0&2&1&0&0\\
2&0&1&6&1&0\\
3&0&0&1&2&0\\
4&0&0&0&1&1
\end{array}
\]
\end{theorem}
\begin{proof}
\begin{enumerate}[label={$\langle1\rangle$\arabic*.}]
\item All maps needed for the two nontrivial rows can be evaluated at
the cusp. We give their matrices and the reason evaluation
determines their global ranks. In the algebraic splitting of
\cref{night:P1.15} write the three generators of $F$ as $(f_0,f_1,f_2)$ of
weights $(0,1,2)$, and those of $G$ as $(g_0,g_1,g_2)$ of
weights $(0,-1,-2)$. The residue of the quotient connection is
\[
R:F_0\longrightarrow G_0,\qquad
R=\begin{pmatrix}0&A&B\\C&0&D_0\\E&D_0&0\end{pmatrix},
\]
where
\begin{align*}
A&=-125(192z-3125)/12288,&
B&=-78125(96z+3125)/393216,\\
C&=-125/512,&E&=78125/32768,\\
D_0&=-146484375/2097152.
\end{align*}
Here the lifted first vector is $w_0+(4/3)ze$.
These entries result from the exact Gauss--Manin reduction of
\cref{night:P1.15}, with the quotient projection taken after lifting $F$.
The full projected residue and both connecting determinants are
recomputed in \href{https://github.com/KokunoYumeto/s6-complex-structure-workbench/blob/83d76022bcd442278b760db9ff438afceed6c6ea/pro_returns/S6_SERIES_GRIND_checkpoint_10_20261120T0211Z/code/current_audits/audit_corrections/code/independent_V_analytic_06.py}{\texttt{independent\_V\_analytic\_06.py}}.
Its JSON output also records the root-line cohomology sums used below.

\item The connecting map of \eqref{eq:P1Alogcotangent}, on relative
cohomology, sends $f\otimes\eta$ to $R(f)\wedge\eta$ at the
cusp. To see the extension class, trivialize logarithmic
forms on the toric covering. A deck translation changes a
relative logarithmic form by its period derivative times
$dq/q$. The constant Cech class of this derivative is exactly
the residue of the quotient connection. This proves the
formula with the indicated frames. It also proves its exterior
power version, with the usual alternating signs.

\item The floor formula gives the following absolute cohomology:
\[
\begin{array}{c|ccccc}
&h^0&h^1&h^2&h^3&h^4\\\hline
h^*F&1&3&2&0&0\\
h^*L^{11}&0&2&3&1&0\\
h^*(\bigwedge^2F)&0&2&3&1&0\\
h^*(L^{11}\otimes F)&2&6&7&3&0.
\end{array}
\]
Every group in this display comes from global sections of the
indicated direct image; none comes from $H^1$ of a line on the
coarse base. Every contributing summand has coarse degree zero,
so evaluation at the cusp is injective on its space of sections.
For the equal-dimensional source and target spaces used below
it is an isomorphism onto the displayed weight subspaces.
Thus fibrewise evaluation computes the ranks of the
absolute connecting maps.

\item For one-forms, the three middle source classes are $f_i\otimes g_i$.
In the ordered target $(g_0g_1,g_0g_2,g_1g_2)$, the matrix is
\[
C_1=\begin{pmatrix}-C&A&0\\-E&0&B\\0&-D_0&D_0\end{pmatrix},
\qquad
\det C_1=
\frac{476837158203125}{281474976710656}(384z+3125).
\]
The degree-zero connecting map is nonzero because $R(f_0)$
has nonzero $g_1,g_2$ coefficients. The degree-two connecting
map, from the classes $f_1g_0g_1,f_2g_0g_2$ to $g_0g_1g_2$,
has rank one because $D_0\ne0$.

\item We must exclude the sole possible zero of $\det C_1$ on the
whole admissible parameter domain. \Cref{night:P1.15} identifies $z=(625\pi i/128)\xi$ on the full complex
parameter plane, and hence $z=-(625\pi/128)c$ when $\xi=ic$.
The unique exceptional parameter is therefore
$\xi=5i/(3\pi)$. By \cref{night:P1.18}, the admissible half-plane is
$\operatorname{Im}\xi>-h_{b,\min}>1$, whereas $5/(3\pi)<1$.
Thus this zero lies outside the whole admissible parameter domain.
The bound used here is
\cref{night:P1.18}.
Consequently the three connecting ranks are $1,3,1$.
The long exact sequence of \eqref{eq:P1Alogcotangent} gives
\[
h^q(\Omega_M^1(\log W))=(0,1,1,0,0).
\]

\item The residue sequence
\[
0\to\Omega_M^1\to\Omega_M^1(\log W)\to\nu_*\mathcal O_Y\to0
\]
is exact even at self-intersections: in a local branch chart it
records independently the coefficient of each $dx_i/x_i$ on
$x_i=0$. The normalization $Y$ is smooth complete toric, so its
only coherent structure cohomology is $H^0=\mathbb C$ \cite[Theorem 9.2.3, p.~410]{CoxLittleSchenck2011}.
The one-form row follows immediately.

\item For two-forms take the second exterior power of
\eqref{eq:P1Alogcotangent}. Its short exact sequence has left
term $h^*(L^{11}\otimes F)$ and right term
$h^*(\bigwedge^2F)$. In degree two the three source classes
are $f_if_j\otimes g_ig_j$, $i<j$. The target is
$F\otimes g_0g_1g_2$. Its connecting matrix is
\[
C_2=\begin{pmatrix}-D_0&D_0&0\\E&0&-B\\0&-C&A\end{pmatrix},
\qquad \det C_2=-\det C_1.
\]
Thus it is an isomorphism, and
$h^3(\Omega_M^2(\log W))=1$.

\item For completeness the ordinary/logarithmic comparison in this
degree is the exact residue resolution
\[
0\to\Omega_M^2\to\Omega_M^2(\log W)
 \to\nu_*\Omega_Y^1(\log\partial Y)
 \to\nu^{[2]}_*\mathcal O_{W^{[2]}}\to0.
\]
The last space is the disjoint normalization of the seven
double surfaces. Local exactness is verified by writing a
logarithmic two-form as a sum of
$a_{ij}\,dx_i/x_i\wedge dx_j/x_j$ and terms with ordinary
differentials. The two residues of an $ij$ term have opposite
sign; their cancellation is the sole compatibility. Every local
double-surface coefficient lifts to one branch, which proves
surjectivity at the last term.
Both final sheaves have zero positive cohomology: the first
is the pushforward of the trivial rank-three toric logarithmic
bundle, and the second has complete toric surface components.
Writing $Q=\Omega_M^2(\log W)/\Omega_M^2$, we obtain
$H^q(Q)=0$ for $q\ge2$. Hence
$h^3(\Omega_M^2)=1$.
In degree zero the last residue map is zero: the fourteen toric
boundary divisors of $Y$ occur in seven opposite pairs; their
primitive normal vectors have opposite signs and their gluing
is a torus translation. This gives the zero $7$-by-$3$ residue
matrix on global logarithmic one-forms. Thus $H^0(Q)=\mathbb C^3$.
The two global logarithmic two-forms come from the
$L^{12},L^{13}$ summands of $L^{11}\otimes F$; their residues
are the independent $f_1,f_2$ directions at the cusp.
Therefore $h^0(\Omega_M^2)=0$.

\item Serre duality \cite[Theorem 4, p.~22]{Serre1955Duality} on this compact smooth complex fourfold gives
$h^{2,1}=h^{2,3}=1$. Finally the analytic Riemann--Roch theorem
\cite[Theorem 25.3.1, p.~187; \S25.4(a), pp.~188--190]{HirzebruchTopologicalMethods} applies to the Dolbeault elliptic
operator of $\Omega_M^p$. No Kähler hypothesis occurs there.
Using $c_1=c_3=0$, the Chern-root product gives
\[
\chi_y(M)=
\frac{(1+y)^4}{240}\,c_2^2+
\left(\frac{(-1+10y-y^2)^2}{144}
       -\frac{(1+y)^4}{120}\right)c_4
=-y+4y^2-y^3.
\]
The exact Chern-number substitution and the resulting polynomial are
verified in \href{https://github.com/KokunoYumeto/s6-complex-structure-workbench/blob/83d76022bcd442278b760db9ff438afceed6c6ea/pro_returns/S6_SERIES_GRIND_checkpoint_10_20261120T0211Z/code/current_audits/audit_corrections/code/independent_V_analytic_06.py}{\texttt{independent\_V\_analytic\_06.py}}. It gives
$\chi(\Omega_M^2)=4$, hence $h^{2,2}=6$.
\Cref{night:P1A.1} and Serre duality supply all other entries of the table.

\end{enumerate}
\end{proof}




\begin{theorem}[Tangent cohomology and the Kuranishi germ]\hypertarget{tag-02R3}{}\label{night:P1A.3}
\StatementRecord{night:P1A.3}{proved}{02R3}{\SeriesRef{night:P1.9}, \SeriesRef{night:P1.10}, \SeriesRef{night:P1.15}, \SeriesRef{night:P1A.1}, \SeriesRef{night:P1A.2}}
For every admissible fourfold,
\[
 (h^0(T_M),h^1(T_M),h^2(T_M),h^3(T_M),h^4(T_M))
 =(1,1,1,2,0).
\]
The complex parameter family has a nonzero Kodaira--Spencer map
at every member and supplies a smooth one-dimensional miniversal
Kuranishi germ.
\end{theorem}
\begin{proof}
\begin{enumerate}[label={$\langle1\rangle$\arabic*.}]
\item Put $a=-256/3125$, and denote the whole reduced cusp fibre by
$W$.  The logarithmic sequence established in the proof of
\cref{night:P1A.2}, tensored with $K_M=h^*L^{-6}$ from
\cref{night:P1A.1}, is
\[
0\longrightarrow h^*L^5
\longrightarrow\Omega^1_M(\log W)\otimes K_M
\longrightarrow h^*(F\otimes L^{-6})\longrightarrow0.
\]
Here $K_{\widehat B}(a)=L^{11}$.  More precisely, in the meromorphic
$L$-frames of \cref{night:P1.15},
$\kappa=d\lambda/(\lambda-a)$ is $\lambda^9$ times the
$L^{11}$-frame: both sides have divisor $3D_0-D_\infty$.
The logarithmic sequence is exact in each chart
$q=z_0\cdots z_r$, where its first arrow sends $dq/q$ to
$\sum_{j=0}^r dz_j/z_j$; the quotient is the relative translation
cotangent bundle $h^*F$.  These arrows agree on the monomial and deck transitions and descend through the free finite
root-chart actions.

\item For a summand $L^n$, the coarse degree is
$d(n)=\lfloor4n/5\rfloor+\lfloor-3n/4\rfloor$.
Projection formula and
$R^jh_*\mathcal O_M=\bigwedge^jG$ compute the two end terms.
For $h^*L^5$ the weights in relative degrees $j=0,1,2,3$ are
\[
\{5\},\quad\{5,4,3\},\quad\{4,3,2\},\quad\{2\}.
\]
Weights $4,5$ have degree zero and weights $2,3$ have degree $-1$.
Thus $h^q(h^*L^5)=(1,2,1,0,0)$.
The three weights of $F\otimes L^{-6}$ are $-4,-5,-6$.
After tensoring with $\bigwedge^jG$, every weight is between
$-9$ and $-4$.  We have $d(-9)=-2$ and $d(n)=-1$ for
$-8\le n\le-4$.  Weight $-9$ occurs once in each of relative
degrees $2$ and $3$, so
\[
h^q\bigl(h^*(F\otimes L^{-6})\bigr)=(0,0,0,1,1).
\]
The base has cohomological dimension one, so the two-column Leray
spectral sequences have no further differential.  Every connecting
map for the displayed logarithmic sequence has zero target.
Consequently
\[
h^q\bigl(\Omega^1_M(\log W)\otimes K_M\bigr)=(1,2,1,1,1).
\]

\item Let $\nu:Y\to W$ be the branch normalization.  The residue
sequence, with its branchwise maps, is
\[
0\longrightarrow\Omega^1_M\otimes K_M
\longrightarrow\Omega^1_M(\log W)\otimes K_M
\longrightarrow\nu_*\mathcal O_Y\otimes K_M|_W
\longrightarrow0.
\]
Since $K_M$ is pulled back from the base, its restriction to $W$ is
a constant line.  The complete smooth toric variety $Y$ has
$H^0(\mathcal O_Y)=\mathbb C$ and $H^{>0}(\mathcal O_Y)=0$
\cite[Theorem~9.2.3, p.~410, with divisor zero]{CoxLittleSchenck2011}.
The sole global section of the logarithmic middle term comes from
$h^*L^5$.  Its meromorphic expression is
$\lambda^{-5}\kappa$ tensored with the $L^{-6}$-frame: the
unique $L^5$ section is $\lambda^4$ times its frame, and
$\kappa$ contributes $\lambda^9$.  Its residue at $a$ is
$a^{-5}\ne0$.  Hence the map on global sections is an isomorphism
onto the residue line, and
\[
h^q(\Omega^1_M\otimes K_M)=(0,2,1,1,1).
\]
The pointwise perfect evaluation pairing
$T_M\otimes(\Omega^1_M\otimes K_M)\to K_M$ induces analytic
Serre duality
\cite[\S13, Theorem~4, p.~22; \S15, p.~23]{Serre1955Duality}.
It identifies $H^q(T_M)^\vee$ with
$H^{4-q}(\Omega^1_M\otimes K_M)$ and gives the stated tangent
cohomology dimensions.

\item We check the class of the parameter family on the
smooth fibre $A=h^{-1}(1)$.  In the fixed frame its period plane
has basis
\[
w_2,\quad w_1,\quad w_0+\frac{625\pi i}{96}\xi e.
\]
Differentiation in $\xi$, followed by projection to $G=H/F$,
gives the rank-one map $P:F\to G$ with sole nonzero matrix entry
$P_{13}=625\pi i/96$ in the ordered bases above and
$(e,w_{-1},w_{-2})$ of the quotient.
Differentiation in the base coordinate gives, by projecting the
connection in \cref{night:P1.15}, the matrix
\[
Q=\frac1{3125(1-a)}
\begin{pmatrix}
-78125/6+300r&-390625/48-375r&0\\
120&0&125/2\\
0&120&50
\end{pmatrix},
\qquad r=\frac{625\pi}{96}(-i\xi).
\]
Its determinant is $3/[625(1-a)^3]\ne0$.
The translation-valued \v Cech cocycles obtained by differentiating
the lattice gluing identify
$H^1(T_A)=\operatorname{Hom}(F,G)$; contraction with a relative
one-form sends the two cocycles to the displayed projected period
derivatives, in the same convention.

\item The normal tangent sequence on this fibre is
\[
0\longrightarrow T_A\longrightarrow T_M|_A
\longrightarrow\mathcal O_A\otimes T_{\widehat B,1}
\longrightarrow0.
\]
Its connecting map sends a base tangent vector to its
Kodaira--Spencer class $Q$.  As $H^0(\mathcal O_A)=\mathbb C$,
the kernel of $H^1(T_A)\to H^1(T_M|_A)$ is exactly
$\mathbb C Q$.  The nonzero matrix $P$ has determinant zero;
every nonzero scalar multiple of $Q$ has nonzero determinant.
Therefore the image of $P$ in $H^1(T_M|_A)$ is nonzero.
The global parameter Kodaira--Spencer cocycle restricts to this
image, so its class in $H^1(T_M)$ is nonzero.
The finite affine actions are fixed, and the cusp multipliers
$\exp(2\pi iC(q,\xi)\nu)$ vary holomorphically in $\xi$.
The same finite fan charts, with $C$ bounded on a small parameter
neighbourhood, give the proper local family of
\cref{night:P1.9}, \cref{night:P1.10}.  Thus the cocycle just computed is
the derivative of this compact family.

\item Douady's versality and semiuniversality theorems
\cite[\S VII.8, pp.~598--599; \S VIII.1, pp.~599--600]{Douady1974}
give a miniversal germ with tangent space $H^1(T_M)=\mathbb C$.
Its local ring has a minimal presentation $\mathbb C\{u\}/I$,
where $I\subset(u)^2$.  The parameter disc classifies by
$u\mapsto g(s)$ with $g'(0)\ne0$, by the preceding step.
Substitution $\mathbb C\{u\}\to\mathbb C\{s\}$ is therefore an
isomorphism of convergent power-series rings.  Since this map factors
through $\mathbb C\{u\}/I$, it annihilates $I$; its injectivity
gives $I=0$.  The miniversal germ is reduced and smooth of dimension
one, and the classifying map is a local isomorphism.  This proves
miniversality of the given family at each member.
\end{enumerate}
\end{proof}




\begin{theorem}[Algebraic dimension and automorphisms]\hypertarget{tag-02R4}{}\label{night:P1A.4}
\StatementRecord{night:P1A.4}{proved}{02R4}{\SeriesRef{night:P1.1}, \SeriesRef{night:P1.9}, \SeriesRef{night:P1.15}, \SeriesRef{night:P1A.1}}
The fourfold has algebraic dimension one and
\[
\Aut^0(M)=\mathbb C,\qquad
\Aut(M)=\mathbb C\times\mathbb Z/2.
\]
In a fixed holomorphic splitting, two members of the full
admissible complex parameter half-plane are biholomorphic
exactly when their parameters are equal.
\end{theorem}
\begin{proof}
\begin{enumerate}[label={$\langle1\rangle$\arabic*.}]
\item An effective divisor dominating the base would restrict to
a nonempty effective divisor on a general complex torus.
Pairing its Chern class with the square of a Kähler form
on that torus is strictly positive. The restriction class is an alternating integral form $J$ on $\Lambda$
fixed by both monodromies:
\[
M_1^tJM_1=J,\qquad M_2^tJM_2=J.
\]
The exact linear system in its fifteen alternating entries has rank
fifteen, so $J=0$.  This contradicts the positive pairing, and every
divisor is vertical.
The pole divisor of a meromorphic function is vertical;
its restriction to a general compact fibre is holomorphic
and therefore constant. The compact reduced graph of the function maps properly to
$\mathbb P^1\times\mathbb P^1$ by $h$ on its first factor.
Remmert's proper-image theorem, in
\cite[\S7, Satz~1, p.~60]{Grauert60}, makes its image analytic.
It is generically a single graph and therefore defines the descended
meromorphic function on the normal base curve, including the three
special values.
Conversely pullbacks from the base occur. Hence the
meromorphic function field is $\mathbb C(t)$ and $a(M)=1$.
Every automorphism consequently induces a Möbius map on
the base. The two distinct multiplicities and the cusp
identify all three values separately, so this map is the
identity.

\item On smooth fibres an automorphism is affine. Its integral
linear part is locally constant and commutes with $M_1,M_2$.
The two simultaneous commutator equations have rational solution
space $\mathbb QI$, as recomputed in
\href{https://github.com/KokunoYumeto/s6-complex-structure-workbench/blob/83d76022bcd442278b760db9ff438afceed6c6ea/pro_returns/S6_SERIES_GRIND_checkpoint_10_20261120T0211Z/code/current_audits/audit_corrections/code/independent_V_extension_06.py}{\texttt{independent\_V\_extension\_06.py}}; their integral invertible
elements are therefore $\pm I$. At the order-five filling an affine lift with
linear part $\epsilon I$ would have to conjugate the same
finite affine generator. For its translation $b$ this
requires
\[
(I-M_1)b=(1-\epsilon)v_1/5\pmod\Lambda.
\]
Applying the invariant integral covector $\ell$ makes the
left side zero. For $\epsilon=-1$ the right side is
$2/5\pmod\mathbb Z$, which is nonzero. Thus every
automorphism is a translation.

\item Let $\mathcal A$ be the sheaf of translations extending to
the cusp model and finite root charts. There is
an exact analytic sheaf sequence
\[
0\to j_*\Lambda\to E_0\to\mathcal A\to0
\]
on the root orbifold. Here $j$ omits the cusp and the stalk
of $j_*\Lambda$ there is $\ker N$.
On a finite root cover the torus evaluation sequence is exact and
equivariant for the prescribed affine action; its translation part
cancels under conjugation.  We descend this sequence on the orbifold, with its finite stabilizer cohomology.

To check the last surjectivity at the cusp, lift a translation
to the toroidal cover. It preserves the fan precisely when
its valuation is an integral translation of the periodic
fan. Subtract the corresponding period; the remaining
characters are nonvanishing holomorphic functions on a
disc and possess logarithms. These give a local section
of $E_0$. The kernel consists exactly of the invariant
integral periods. This also proves that the sequence
accounts for all extendable translations, including
possible component permutations.

\item Write $A=M_1$, $B=M_2$ and
$N_1=I+A+\cdots+A^4$, $N_2=I+B+B^2+B^3$.
The orbifold fundamental group away from the cusp has generators
$g_1,g_2$ with relations $g_1^5=g_2^4=1$.  A lattice cocycle is
therefore a pair $(x_1,x_2)$ with $N_1x_1=N_2x_2=0$.
The degree-one exact sequence for $j_*$ says that it extends across
the cusp precisely when its value on $g_1g_2$ is a coboundary:
\[
x_1+Ax_2=(AB-I)u,\qquad u\in\Lambda.
\]
Subtract the coboundary of $u$.  The resulting pair is $(-Ay,y)$,
where $N_1y=N_2y=0$.  The ambiguity in $u$ is
$t\in\ker(AB-I)$, which changes $y$ by $(B-I)t$.
Thus the cohomology group, with the finite-root torsion, is
\[
H^1(\widehat B,j_*\Lambda)=
\frac{\ker N_1\cap\ker N_2}{(B-I)\ker(AB-I)}.
\]
Primitive integral column bases for its numerator and for
$\ker(AB-I)$ are respectively
\[
K=\begin{pmatrix}
0&0&0\\0&1&0\\1&0&0\\0&0&0\\0&-1&0\\0&0&1
\end{pmatrix},\qquad
T=\begin{pmatrix}
0&0&0\\0&1&0\\1&0&0\\0&0&-1\\0&0&1\\0&0&-1
\end{pmatrix}.
\]
Direct multiplication gives
\[
(B-I)T=KZ,\qquad
Z=\begin{pmatrix}-2&-1&1\\0&-1&1\\1&0&1\end{pmatrix}.
\]
The determinant of $Z$ has absolute value two; an entry and a
$2$-by-$2$ minor have absolute value one.  Its Smith factors are
$1,1,2$, proving $H^1(\widehat B,j_*\Lambda)=\mathbb Z/2$.
The norm-kernel ranks, primitive column bases, full identity
$(B-I)T=KZ$ and Smith factors are recomputed exactly in
\href{https://github.com/KokunoYumeto/s6-complex-structure-workbench/blob/83d76022bcd442278b760db9ff438afceed6c6ea/pro_returns/S6_SERIES_GRIND_checkpoint_10_20261120T0211Z/code/current_audits/audit_corrections/code/independent_V_analytic_06.py}{\texttt{independent\_V\_analytic\_06.py}}.
\item Now $H^0(j_*\Lambda)=0$ by \cref{night:P1.1},
$H^0(E_0)=\mathbb C$, and $H^1(E_0)=0$ by the floor
calculation. The exact sequence gives
\[
0\to\mathbb C\to H^0(\mathcal A)\to\mathbb Z/2\to0.
\]
The half-period
$\tau=(0,1,1,1,1,1)^t/2$ is fixed modulo $\Lambda$
by every monodromy matrix. It extends over the finite
quotients because the translations commute there.
At the cusp $N\tau$ is integral; primitiveness of $N$
provides $\lambda\in\Lambda$ with $N\lambda=N\tau$.
The representative $\tau-\lambda\in\ker N$ is a bounded
angular half-period and extends on every toric chart.
It is a nonzero order-two global section. Since the
additive group $\mathbb C$ has no nonzero torsion, it
splits the displayed exact sequence. Translations
commute, giving the asserted full group.

\item For an isomorphism between two members, the same intrinsic
base argument applies. Its integral linear part still
lies in the common centralizer. Either scalar $\pm I$
preserves a complex period plane exactly by equality
of the two planes. The graphs
$(\sigma+\xi\alpha e)(P)$ are equal precisely when their
parameters $\xi$ agree, since $e\notin F$ and $\alpha$
is nonzero. Translation does not change this condition.
This proves the parameter assertion.

\end{enumerate}
\end{proof}




\begin{theorem}[No antiholomorphic isomorphisms]\hypertarget{tag-02QQ}{}\label{night:P1.17}
\StatementRecord{night:P1.17}{proved}{02QQ}{\SeriesRef{night:P1.1}}
For the cusp marking $S$ fixed in \cref{night:P1.1} and the corresponding Kuhn fan,
none of the completed fourfolds in this family admits an
anti-holomorphic automorphism.  In particular none admits a real
structure.  More strongly, two members with this same cusp marking
are never anti-biholomorphic.
\end{theorem}
\begin{proof}
\begin{enumerate}[label={$\langle1\rangle$\arabic*.}]
\item We first recover the base intrinsically.  An effective divisor on
$M_c$ cannot be horizontal: its first Chern class vanishes by
$H^2(M_c;\mathbb Z)=0$, whereas its nonzero effective restriction to a
general torus fibre would have positive pairing with the square of a
flat K\"ahler class.  The pole divisor of any meromorphic function is
therefore vertical.  On every regular fibre outside that divisor the
function is holomorphic and hence constant.  It factors through a
meromorphic function on the base, which is rational on
$\mathbb P^1$.  Thus the function field is exactly $\mathbb C(\lambda)$.
An anti-biholomorphism consequently induces an anti-M\"obius map of
the base.  The two distinct multiplicities $5,4$ and the singular cusp
support distinguish the three special values, so this base map fixes
$\infty,0,\lambda_*$.  As those points are real, the map is precisely
$\lambda\mapsto\bar\lambda$.

\item A holomorphic map between compact complex tori is affine: after
lifting to their vector spaces, its derivative is periodic and
therefore constant, since it descends to holomorphic functions on a
compact torus.  Applying this to an anti-holomorphic map after
conjugating its source gives the same assertion for its anti-linear
lift.  Its induced integral map on $H_1$ is constant along the
regular base.  Complex conjugation reverses the two finite meridians,
so this matrix must simultaneously conjugate $M_i$ to $M_i^{-1}$.
Solving these exact integral linear equations gives only $R$ and
$-R$ among unimodular matrices, where
\[
R=\begin{pmatrix}
1&0&0&0&0&0\\
2&0&0&0&-1&0\\
-1&0&-1&0&0&0\\
-1&0&0&-1&0&0\\
2&-1&0&0&0&0\\
-1&1&1&1&1&1
\end{pmatrix},\qquad R^2=I.
\]
Indeed the simultaneous intertwiner space over $\mathbb Q$ has
dimension one and this primitive matrix is unimodular; determinant
one in absolute value permits only the scalar factors $\pm1$.
The exact anti-intertwiner equations, their one-dimensional rational
solution space, unimodularity of $R$, the based cusp block and the
specific ray obstruction below are recomputed in
\href{https://github.com/KokunoYumeto/s6-complex-structure-workbench/blob/83d76022bcd442278b760db9ff438afceed6c6ea/pro_returns/S6_SERIES_GRIND_checkpoint_10_20261120T0211Z/code/current_audits/audit_corrections/code/independent_V_markings_06.py}{\texttt{independent\_V\_markings\_06.py}} and its JSON output.

\item The finite affine cap excludes $G=R$.  Normalize the pullback by
$t=u^{m_i}$; any base-preserving isomorphism lifts to this canonical
root cover and conjugates its deck generator to the inverse.
At its central torus the translation part $t_i$ would satisfy
\[
(I-M_i^{-1})t_i=-(v_i+Gv_i)/m_i\pmod\Lambda.
\]
Applying $\ell$, and using $\ell R=\ell$, gives
$0=-2\ell(v_i)/m_i$ in $\mathbb R/\mathbb Z$, impossible already
for $m_i=5$.  Hence $G=-R$ is necessary.

\item We now use the cusp marking, keeping track of the conjugation of based meridians.  Conjugation sends
$g_0=(g_1g_2)^{-1}$ to $g_2g_1=g_2g_0^{-1}g_2^{-1}$.
Thus the matrix in a fixed local cusp frame is
$H=M_2^{-1}G$, and it anti-conjugates $M_0$.
Exact multiplication in the marking $S$ gives
\[
S^{-1}HS=
\begin{pmatrix}
1&0&0&-1&-1&0\\
0&1&0&-2&0&0\\
1&-1&-1&1&-1&-1\\
0&0&0&-1&0&0\\
0&0&0&0&-1&0\\
0&0&0&-1&1&1
\end{pmatrix}
=\begin{pmatrix}P&Q\\0&-P\end{pmatrix},\quad
P=\begin{pmatrix}1&0&0\\0&1&0\\1&-1&-1\end{pmatrix}.
\tag{1.17.1}
\]
Changing the based cusp path by a cusp loop multiplies by a power of
$M_0$ and does not change $P$.

\item We explain why an arbitrary translation cannot remove the obstruction
in (1.17.1).  In the additive torus coordinates whose first periods
are the $a_j$, the anti-linear map must have the form
$\zeta' =P\bar\zeta+t(q)$, because its values on these three real
periods already determine its linear part.  In multiplicative
coordinates it is
$z'=u(q)\bar z^{-P}$.
The source zero section limits to a smooth point of the central fibre.
An isomorphism takes it to another smooth point.  The central fibre
is irreducible, and its smooth locus is exactly the dense torus of
its toric normalization: every toric boundary divisor is a
normal-crossings self-identification.  In a chart about this image,
the translation therefore has the form $q^v$ times nonvanishing
holomorphic units, with $v\in\mathbb Z^3$.  Translating by the period
$v$ changes the selected component chart; the units act by torus
translation.  Neither changes the linear action $-P$ on the fan.

\item The primitive rays of the normal fan of that component are
\[
\{\pm v:0\ne v\in\{0,1\}^3\}.
\]
An extending torus monomial isomorphism must permute these rays.
This can be seen directly on a one-parameter arc $z=t^v$ meeting the
generic point of its boundary divisor: if the image vector lies in
the interior of a cone of dimension at least two, that divisor is
contracted, which is impossible for an isomorphism.  Complex
conjugation itself preserves the toric charts, so the same criterion
applies to the anti-monomial map above.
But $-Pe_2=(0,-1,1)^t$ is not one of the displayed rays.
This is a contradiction.  We have exhausted both possible integral
linear parts, so no anti-biholomorphism exists.


\end{enumerate}
\end{proof}


\section{The boundary of the family}

Fix one splitting and write the complex parameter as $\xi$, so that
$F_\xi=(\sigma+\xi\alpha e)(P)$. Put
$c_*=-\min h$ and $\mathcal H=\{\operatorname{Im}\xi>c_*\}$.
The real family is $\xi=ic$, $c\in\mathbb R$.
By an analytic extension across a boundary point we mean a proper
Hausdorff complex-analytic family on a full complex disc whose
restriction to the nonempty open part lying in $\mathcal H$ is
the prescribed family. At infinity we use the coordinate $u=1/\xi$.
This definition permits singular central spaces, multiplicities,
and additional central components.

The stability of complex tori under deformation in the large is
established in \cite[definition, p.~261; Theorem 2.1, p.~264]{Catanese2004}.
The relative-form calculation below writes the torus quotient
on the regular analytic open set used in the boundary argument.

\begin{theorem}[No analytic extension across the boundary]\hypertarget{tag-02R5}{}\label{night:P1A.5}
\StatementRecord{night:P1A.5}{proved}{02R5}{\SeriesRef{night:P2.4}, \SeriesRef{night:P1A.1}, \SeriesRef{night:P1.8}; \cite[Theorem 1.6, preprint p.~5]{BM97}; \cite[\S6, Hauptsatz I; \S7, Satz 1 and Satz 5]{Grauert60}}
There is no such analytic extension across any finite boundary
point of $\mathcal H$, or across infinity. The same conclusion
holds after any finite ramified change of the local boundary
coordinate. In particular neither end of the real interval
$c_*<c<+\infty$ has a proper Hausdorff analytic extension in
this sense.
\end{theorem}
\begin{proof}
\begin{enumerate}[label={$\langle1\rangle$\arabic*.}]
\item We first identify the moving pencil on each known member.
By \cref{night:P1A.1}, $K_M^{-4}=h^*L^{24}$ and the fibres of $h$
are connected.  For the coarse map
$\rho:\widehat B\to\mathbb P^1$, the root-divisor formula gives
\[
 \rho_*L^{24}
 =\mathcal O_{\mathbb P^1}
   \bigl(\lfloor96/5\rfloor\,\infty+
                \lfloor-72/4\rfloor\,0\bigr)
 =\mathcal O_{\mathbb P^1}(19\infty-18\,0).
\]
Its sections are spanned by $\lambda^{18},\lambda^{19}$.
On a root chart the invariant sections have a common local
monomial factor, and their ratio cancels that factor.  Thus
$h^0(M,K_M^{-4})=2$, and the moving part of this complete
system is exactly the marked base map, up to a projective
coordinate change.  In particular its base ideal is invertible
on every member of the prescribed family.

\item Suppose that a proper Hausdorff analytic extension over a full
parameter disc exists.  Take the reduced component containing
the given nonempty open family and resolve it.  Analytic
resolution applies on a relatively compact neighbourhood of
the compact central fibre and can be kept unchanged over the
given smooth locus
\cite[Theorem 1.6 and its following analytic-neighbourhood
paragraph, preprint p.~5]{BM97}.
Write $p:Z\to\Delta$ for the resulting proper map from a smooth
irreducible space.  The proper images of the exceptional and
critical loci are analytic
\cite[\S7, Satz 1, p.~60]{Grauert60}; they omit the given open
parameter set.  After shrinking the disc they are supported
at its centre.  Hence $p$ is a proper holomorphic submersion
over $\Delta^*$, with compact fourfold fibres.

Put $\mathcal J=\omega_{Z/\Delta}^{-4}$.
The sheaves $R^jp_*\mathcal J$ are coherent
\cite[\S6, Hauptsatz I, p.~59]{Grauert60}.
Multiplication by a nonzero parameter germ is injective on
$\mathcal J$ and on $p_*\mathcal J$, so $p_*\mathcal J$ is locally free
over the disc and free after shrinking.
The exact sequence for a fibre over $t$ gives
\[
 0\longrightarrow(p_*\mathcal J)_t/\mathfrak m_t(p_*\mathcal J)_t
 \longrightarrow H^0(Z_t,\mathcal J_t)
 \longrightarrow
 \ker\bigl((t'-t): (R^1p_*\mathcal J)_t
                       \longrightarrow(R^1p_*\mathcal J)_t\bigr)
 \longrightarrow0.
\]
The last term is supported in the torsion support of the
coherent sheaf $R^1p_*\mathcal J$, a closed discrete analytic
set.  At a point of the given open family outside this support,
the displayed comparison shows that $p_*\mathcal J$ has rank
two.  Shrinking makes the support empty away from the centre.
Thus a frame of $p_*\mathcal J$ gives the complete two-dimensional
linear system on each punctured fibre.  

\item The two evaluation sections define a coherent base ideal and
a meromorphic map $Z\dashrightarrow\mathbb P^1$ over the full
disc.  The ideal is locally invertible on the given open family
by the first step.  Its non-invertibility locus is analytic;
its proper image in $\Delta$ is therefore discrete and can be
removed away from the centre.  Blowing up that ideal and
resolving the graph consequently changes only the central
fibre.  It gives a proper holomorphic map to
$\mathbb P^1\times\Delta$.  The construction uses the graph
of the coherent evaluation map, so its coefficients have
finite meromorphic order in the parameter.

Its critical image is proper analytic.  Every component
dominating the parameter disc meets the given open family,
where there are precisely three critical values.  Their
degrees over the parameter are one.  The multiplicities five
and four distinguish two of them, and the third is the cusp.
After shrinking, their normalizations give three distinct
labelled meromorphic sections, and all vertical exceptional
components lie over the centre.  The projective transformation
taking these sections to $\infty,0,\lambda_*$ has meromorphic
coefficients and is invertible on $\Delta^*$.  Resolving its
graph gives the marked map $H$ over the disc, still
unchanged on punctured fibres.  Its regular restriction is a
proper holomorphic submersion
\[
 q:W\longrightarrow
 B=\bigl(\mathbb P^1\setminus\{0,\infty,\lambda_*\}\bigr)
          \times\Delta^*.
\]
Ehresmann transport makes the number of components of a fibre
locally constant.  The base is connected and the known fibres
are connected, so every fibre of $q$ is connected.

\item We need the torus structure on an analytically open dense
part of this regular family.  The coherent sheaf
$\mathcal E=q_*\Omega^1_{W/B}$ has generic rank three.
Generic cohomology and base change for this flat bundle gives
a proper analytic exceptional set outside which $\mathcal E$
is locally free of rank three and its fibre is the space of
holomorphic one-forms.  The constant-dimension comparison is
\cite[\S7, Satz 5, pp.~63--64]{Grauert60}.
Use the global coherent evaluation map
\[
 q^*\mathcal E\longrightarrow\Omega^1_{W/B}.
\]
Its cokernel has analytic support.  The proper image of that
support is analytic in $B$ and omits the nonempty part of the
known torus open set outside the generic exceptional set.
Enlarge the exceptional set by this image.  On its complement
evaluation is an isomorphism of rank-three bundles.

Relative exterior differentiation defines an
$\mathcal O_B$-linear map
$\mathcal E\to q_*\Omega^2_{W/B}$.  It vanishes on the known
torus open set.  The target is torsion-free, since it is the
proper direct image of a vector bundle under a submersion.
The identity theorem therefore makes this map zero.  Each
remaining compact connected fibre has three everywhere
independent closed holomorphic one-forms.  Their dual vector
fields commute and are complete by compactness.  Their flows
give a transitive holomorphic $\mathbb C^3$-action: every orbit
is open and the fibre is connected.  Its stabilizer $\Gamma$
is discrete, and compactness of $\mathbb C^3/\Gamma$ makes
$\Gamma$ a full real lattice.  This is the complex-torus
structure on each such fibre.

\item Let $E\subset B$ be the proper analytic exceptional set just
removed.  On the universal cover of $B\setminus E$, use
Gauss--Manin transport from the given open family to mark
integral $H^1$.  Integration of the closed relative forms
gives a holomorphic map to $\operatorname{Gr}(3,6)$.  On the
known open set it is the explicit graph
$F_\xi=(\sigma+\xi\alpha e)(P)$.  The identity theorem on the
connected cover identifies these two maps everywhere.
The explicit graph is a rank-three holomorphic plane for
every finite $\xi$, including parameters for which it is not
real-transverse.  At infinity the same assertion is used with
$\xi=1/u$ on the punctured parameter disc.  All changes of
continuation marking are real integral and preserve the
condition $F\cap\bar F=0$.

\item Let $\xi_0$ be a finite boundary point, so
$\operatorname{Im}\xi_0=c_*=-\min h$.
The minimum may occur at a finite root point.  Continuity
there nevertheless gives regular $\lambda$ with $h(\lambda)$
arbitrarily close to $\min h$.  Every sufficiently small full
parameter disc about $\xi_0$ therefore contains noncentral
points of
\[
 \operatorname{Im}\xi+h(\lambda)=0.
\]
This equation defines nonempty smooth real hypersurface
pieces in the complex surface $B$, since its derivative in
$\operatorname{Im}\xi$ is one.  A proper complex analytic
subset of $B$ has real dimension at most two; it cannot
contain those three-dimensional pieces.  Choose such a
point outside $E$.  Its fibre is a complex torus by the fourth
step, so its period plane is real-transverse.  The fifth step
identifies that plane with $F_\xi$, whereas the exact criterion
of \cref{night:P1.8} gives
$F_\xi\cap\overline{F_\xi}\ne0$ at the chosen point.
This is a contradiction.

At infinity fix a regular $\lambda$ and set
\[
 \xi=R-i h(\lambda),\qquad
 u=\frac1{R-i h(\lambda)},\qquad R\in\mathbb R.
\]
These are noncentral points of the same failure locus and
tend to $u=0$ as $|R|\to\infty$.  Varying $\lambda$ and $R$
again gives real hypersurface pieces in every punctured
$u$-disc, so one can choose a point outside $E$ and obtain the
same contradiction.  Under a finite ramified change of the
boundary coordinate, the parameter map has nonzero derivative
off the centre after shrinking.  The inverse images of these
hypersurface pieces remain nonempty real hypersurfaces there,
and the argument is unchanged.

All contradictions occur on a noncentral compact smooth
torus fibre of the recovered marked map.  The reductions and
modifications are proper and preserve the prescribed open
family; they impose no smoothness, multiplicity or component
condition on the original central fibre.  The permitted
central spaces therefore do not avoid the contradiction.
Both nonextension assertions follow.
\end{enumerate}
\end{proof}




\begin{theorem}[Families over a punctured disc]\hypertarget{tag-02R6}{}\label{night:P1A.6}
\StatementRecord{night:P1A.6}{proved}{02R6}{\SeriesRef{night:P1A.3}, \SeriesRef{night:P1A.4}}
Let $p:\mathcal M\to\Delta^*$ be a proper holomorphic submersion
whose fibres are members of this fourfold family. Its classifying
parameter has an interior limit. After a base change of degree at most two, it is
isomorphic to the pullback of the constructed family by
a holomorphic map of the full disc to $\mathcal H$, and
therefore has a proper smooth filling. Such a family
cannot approach either parameter end.
\end{theorem}
\begin{proof}
\begin{enumerate}[label={$\langle1\rangle$\arabic*.}]
\item At every punctured parameter, \cref{night:P1A.3} supplies local
miniversality and a local holomorphic classifying map to
$\mathcal H$. \Cref{night:P1A.4} makes this map unique on each geometric
point. Thus these local maps agree on overlaps and define
one holomorphic function $\xi:\Delta^*\to\mathcal H$.
There is no quotient by parameter translations, by \cref{night:P1A.4}.

\item The Cayley coordinate
\[
 z=\frac{\xi-i(c_*+1)}{\xi-i(c_*-1)}
\]
has $|z|<1$. It extends holomorphically across the puncture
by the removable-singularity theorem. Its value there has
modulus strictly less than one, since a boundary value
would force $z$ to be constant by the maximum-modulus
principle. The inverse Cayley map therefore extends
$\xi$ with value in $\mathcal H$.

\item Local identifications with the pulled-back constructed
family differ by its automorphisms. These form the constant
holomorphic group $\mathbb C\times\mathbb Z/2$ over the
parameter half-plane: the translation field and the
half-period in \cref{night:P1A.4} extend over the entire family.
The finite component torsor on $\Delta^*$ is killed by
the double cover. For a coordinate $z$ on $\Delta^*$, the functions $z$ and $z^{-1}$
confine the holomorphic hull of each compact set to a compact annulus;
$z$ also separates points and supplies a local coordinate. Thus
$\Delta^*$ is Stein. Cartan's Theorem~B
\cite[Expos\'e~19, \S5, pp.~19-14--19-15]{Cartan1952B_AB},
applied to its coherent structure sheaf, gives
$H^1(\Delta^*,\mathcal O)=0$, so the remaining additive
holomorphic torsor is a coboundary. Thus after that cover the local
identifications glue to a global isomorphism. Pulling
back the already constructed proper smooth family along
the extended classifying map supplies the stated filling.

\end{enumerate}
\end{proof}


